\documentclass[12pt, a4paper]{book}

% ── Encodage et langue ──────────────────────────────────────────────
\usepackage[utf8]{inputenc}
\usepackage[T1]{fontenc}
\usepackage[french]{babel}

% ── Mise en page ────────────────────────────────────────────────────
\usepackage[margin=2.5cm]{geometry}
\usepackage{fancyhdr}
\usepackage{titlesec}
\usepackage{tocloft}

% ── Mathématiques ───────────────────────────────────────────────────
\usepackage{amsmath, amssymb, amsthm, mathtools}
\usepackage{bbm}

% ── Code source ─────────────────────────────────────────────────────
\usepackage{listings}
\usepackage{xcolor}

% ── Graphiques et boîtes ────────────────────────────────────────────
\usepackage{tikz}
\usepackage{mdframed}
\usepackage{tcolorbox}
\tcbuselibrary{theorems, skins, breakable}
\usepackage{booktabs}
\usepackage{array}
\usepackage{hyperref}
\usepackage{graphicx}

% ── Couleurs du projet ──────────────────────────────────────────────
\definecolor{suva_blue}{RGB}{30, 100, 180}
\definecolor{suva_dark}{RGB}{15, 20, 40}
\definecolor{code_bg}{RGB}{245, 247, 250}
\definecolor{code_kw}{RGB}{0, 80, 200}
\definecolor{code_str}{RGB}{180, 40, 40}
\definecolor{code_cmt}{RGB}{80, 130, 80}
\definecolor{warn_bg}{RGB}{255, 248, 220}
\definecolor{warn_border}{RGB}{200, 150, 0}
\definecolor{info_bg}{RGB}{230, 245, 255}
\definecolor{info_border}{RGB}{30, 100, 180}
\definecolor{success_bg}{RGB}{230, 255, 235}
\definecolor{success_border}{RGB}{40, 160, 80}

% ── Configuration listings Python ───────────────────────────────────
\lstdefinestyle{python}{
  language=Python,
  basicstyle=\ttfamily\small,
  keywordstyle=\color{code_kw}\bfseries,
  stringstyle=\color{code_str},
  commentstyle=\color{code_cmt}\itshape,
  backgroundcolor=\color{code_bg},
  frame=single, framerule=0.5pt, rulecolor=\color{gray!40},
  numbers=left, numberstyle=\tiny\color{gray},
  stepnumber=1, numbersep=8pt,
  breaklines=true, breakatwhitespace=true,
  showstringspaces=false,
  tabsize=4,
  morekeywords={assert, yield, with, as, lambda, True, False, None},
}

\lstdefinestyle{solidity}{
  language=C++,
  basicstyle=\ttfamily\small,
  keywordstyle=\color{code_kw}\bfseries,
  stringstyle=\color{code_str},
  commentstyle=\color{code_cmt}\itshape,
  backgroundcolor=\color{code_bg},
  frame=single, framerule=0.5pt, rulecolor=\color{gray!40},
  numbers=left, numberstyle=\tiny\color{gray},
  stepnumber=1, numbersep=8pt,
  breaklines=true,
  morekeywords={uint256, address, bytes, bytes32, memory, storage, external,
                view, pure, returns, mapping, struct, event, emit,
                require, revert, modifier, payable, public, private, internal},
}

\lstdefinestyle{bash}{
  language=bash,
  basicstyle=\ttfamily\small,
  keywordstyle=\color{code_kw}\bfseries,
  backgroundcolor=\color{suva_dark!10},
  frame=single, framerule=0.5pt, rulecolor=\color{gray!60},
  breaklines=true,
  showstringspaces=false,
}

% ── Environnements personnalisés ─────────────────────────────────────
\newtcolorbox{theorembox}[2][]{
  colback=info_bg, colframe=info_border,
  fonttitle=\bfseries, title={Théorème~: #2}, #1
}
\newtcolorbox{defbox}[2][]{
  colback=success_bg, colframe=success_border,
  fonttitle=\bfseries, title={Définition~: #2}, #1
}
\newtcolorbox{warningbox}[1][]{
  colback=warn_bg, colframe=warn_border,
  fonttitle=\bfseries, title={Attention}, #1
}
\newtcolorbox{exercice}[2][]{
  colback=white, colframe=suva_blue!60,
  fonttitle=\bfseries, title={Exercice #2}, #1, breakable
}
\newtcolorbox{solution}[1][]{
  colback=gray!5, colframe=gray!40,
  fonttitle=\bfseries\itshape, title={Solution}, #1, breakable
}
\newtcolorbox{importantbox}[1][]{
  colback=suva_blue!10, colframe=suva_blue,
  fonttitle=\bfseries, title={Point clé}, #1
}

% ── Commandes mathématiques ──────────────────────────────────────────
\newcommand{\Zq}{\mathbb{Z}_q}
\newcommand{\Rq}{R_q}
\newcommand{\ZZ}{\mathbb{Z}}
\newcommand{\NN}{\mathbb{N}}
\newcommand{\norm}[1]{\left\|#1\right\|}
\newcommand{\inftynorm}[1]{\left\|#1\right\|_\infty}
\newcommand{\pmod}[1]{\ (\mathrm{mod}\ #1)}
\newcommand{\NTT}{\mathrm{NTT}}
\newcommand{\INTT}{\mathrm{NTT}^{-1}}

% ── En-têtes et pieds de page ────────────────────────────────────────
\pagestyle{fancy}
\fancyhf{}
\fancyhead[LE,RO]{\thepage}
\fancyhead[LO]{\nouppercase{\rightmark}}
\fancyhead[RE]{\nouppercase{\leftmark}}
\renewcommand{\headrulewidth}{0.4pt}

% ── Titre du document ────────────────────────────────────────────────
\title{
  \vspace{-1cm}
  {\Huge\bfseries\color{suva_blue} Guide d'Apprentissage\\[0.3cm]
  \Large Cryptographie Post-Quantique\\
  \normalsize pour le Projet SuVa / ETHDILITHIUM}\\[1cm]
  \large Théorie, Code et Exercices Pratiques
}
\author{Basé sur le codebase ZKNox / SuVa Protocol}
\date{Mars 2026}

% ════════════════════════════════════════════════════════════════════
\begin{document}

\maketitle
\tableofcontents

% ── Préface ──────────────────────────────────────────────────────────
\chapter*{Préface}
\addcontentsline{toc}{chapter}{Préface}

Ce guide est conçu pour te permettre de maîtriser progressivement
toutes les briques nécessaires pour contribuer efficacement au projet
\textbf{SuVa / ETHDILITHIUM}.

\bigskip
\begin{tcolorbox}[colback=suva_blue!10, colframe=suva_blue,
  title={\bfseries Comment utiliser ce guide}, fonttitle=\bfseries]
\begin{enumerate}
  \item Lis la \textbf{théorie} de chaque section.
  \item Tape \textbf{toi-même} les exemples de code (ne copie-colle pas).
  \item Fais les \textbf{exercices} avant de regarder les solutions.
  \item Relis le code source du projet avec le chapitre correspondant ouvert.
\end{enumerate}
\end{tcolorbox}

\bigskip
\textbf{Prérequis minimum~:} Python 3.10+, notions d'arithmétique de lycée.

\bigskip
\textbf{Ordre de lecture recommandé~:}
\begin{center}
\begin{tikzpicture}[node distance=1.2cm,
  box/.style={draw=suva_blue, rounded corners, fill=info_bg,
               minimum width=4cm, minimum height=0.7cm, font=\small}]
  \node[box] (a) {Ch.1 : Guide de test};
  \node[box, below of=a] (b) {Ch.2 : Arithmétique modulaire};
  \node[box, below of=b] (c) {Ch.3 : Anneaux polynomiaux};
  \node[box, below of=c] (d) {Ch.4 : NTT};
  \node[box, below of=d] (e) {Ch.5 : Réseaux de treillis};
  \node[box, below of=e] (f) {Ch.6 : Dilithium complet};
  \node[box, below of=f] (g) {Ch.7 : Ethereum / Solidity};
  \node[box, below of=g] (h) {Ch.8 : Falcon};
  \foreach \x/\y in {a/b, b/c, c/d, d/e, e/f, f/g, g/h}
    \draw[->, thick, suva_blue] (\x) -- (\y);
\end{tikzpicture}
\end{center}

% ── Inclusion des chapitres ───────────────────────────────────────────
\documentclass[12pt, a4paper]{article}
\usepackage[utf8]{inputenc}
\usepackage[T1]{fontenc}
\usepackage[french]{babel}
\usepackage[margin=2.5cm]{geometry}
\usepackage{amsmath, amssymb}
\usepackage{listings}
\usepackage{xcolor}
\usepackage{tcolorbox}
\tcbuselibrary{theorems, skins, breakable}
\usepackage{hyperref}
\usepackage{booktabs}
\usepackage{fancyhdr}
\usepackage{titlesec}
\usepackage{enumitem}
\usepackage{tikz}
\usetikzlibrary{shapes, arrows.meta, positioning}

% ── Couleurs ────────────────────────────────────────────────
\definecolor{suva_blue}{RGB}{30, 100, 180}
\definecolor{code_bg}{RGB}{245, 247, 250}
\definecolor{code_kw}{RGB}{0, 80, 200}
\definecolor{code_str}{RGB}{180, 40, 40}
\definecolor{code_cmt}{RGB}{80, 130, 80}
\definecolor{warn_bg}{RGB}{255, 248, 220}
\definecolor{warn_border}{RGB}{200, 150, 0}
\definecolor{success_bg}{RGB}{230, 255, 235}
\definecolor{success_border}{RGB}{40, 160, 80}
\definecolor{dark_bg}{RGB}{25, 30, 50}

% ── Listings ─────────────────────────────────────────────────
\lstdefinestyle{bash}{
  basicstyle=\ttfamily\small\color{white},
  backgroundcolor=\color{dark_bg},
  frame=single, framerule=0pt,
  breaklines=true, showstringspaces=false,
  moredelim=[is][\color{green!70}]{\$}{\$},
}
\lstdefinestyle{python}{
  language=Python,
  basicstyle=\ttfamily\small,
  keywordstyle=\color{code_kw}\bfseries,
  stringstyle=\color{code_str},
  commentstyle=\color{code_cmt}\itshape,
  backgroundcolor=\color{code_bg},
  frame=single, rulecolor=\color{gray!40},
  numbers=left, numberstyle=\tiny\color{gray},
  breaklines=true, tabsize=4,
  showstringspaces=false,
}
\lstdefinestyle{output}{
  basicstyle=\ttfamily\small\color{black},
  backgroundcolor=\color{gray!10},
  frame=single, rulecolor=\color{gray!50},
  breaklines=true,
}

% ── Boîtes ───────────────────────────────────────────────────
\newtcolorbox{cmdbox}[1]{
  colback=dark_bg, colframe=suva_blue, colupper=white,
  fonttitle=\ttfamily\bfseries\color{white}, title={\texttt{\$} #1}
}
\newtcolorbox{infobox}[1][]{
  colback=blue!5, colframe=suva_blue,
  fonttitle=\bfseries, title={Info}, #1
}
\newtcolorbox{warningbox}[1][]{
  colback=warn_bg, colframe=warn_border,
  fonttitle=\bfseries, title={Attention}, #1
}
\newtcolorbox{successbox}[1][]{
  colback=success_bg, colframe=success_border,
  fonttitle=\bfseries, title={Succès attendu}, #1
}
\newtcolorbox{errorbox}[1][]{
  colback=red!5, colframe=red!60,
  fonttitle=\bfseries, title={Erreur connue / Solution}, #1
}

% ── En-têtes ─────────────────────────────────────────────────
% ── Commande signature ──────────────────────────────────────────────
\newcommand{\auteur}{Mohamed Lamine MBENGUE}
\newcommand{\specialites}{Sécurité · DevOps · Dev Full Stack Web \& Mobile}
\newcommand{\tel}{+221 78 178 44 36}
\newcommand{\email}{mbenguemohamedlamine2022@gmail.com}
\newcommand{\portfolio}{portfolio-alamine.web.app}

\pagestyle{fancy}
\fancyhf{}
\fancyhead[L]{\color{suva_blue}\textbf{SuVa / ETHDILITHIUM}}
\fancyhead[R]{\color{gray}Guide de Test}
\fancyfoot[L]{\footnotesize\color{gray}\auteur}
\fancyfoot[C]{\thepage}
\fancyfoot[R]{\footnotesize\color{gray}\portfolio}
\renewcommand{\headrulewidth}{0.4pt}
\renewcommand{\footrulewidth}{0.3pt}

% ────────────────────────────────────────────────────────────
\begin{document}

\begin{center}
  {\color{suva_blue}\rule{\linewidth}{2pt}}\\[0.5cm]
  {\Huge\bfseries\color{suva_blue} Guide de Test du Projet}\\[0.3cm]
  {\large\color{gray} SuVa / ETHDILITHIUM --- Mise en Route Complète}\\[0.2cm]
  {\small Mars 2026}\\[0.3cm]
  {\color{suva_blue}\rule{\linewidth}{2pt}}
\end{center}

\begin{tcolorbox}[colback=suva_blue!5, colframe=suva_blue!40,
  left=6pt, right=6pt, top=4pt, bottom=4pt]
\small
\begin{tabular}{ll}
  \textbf{Auteur} & \auteur \\
  \textbf{Spécialités} & \specialites \\
  \textbf{Tél} & \tel \quad \textbf{Email} : \email \\
  \textbf{Portfolio} & \href{https://\portfolio}{\texttt{\portfolio}} \\
\end{tabular}
\end{tcolorbox}

\tableofcontents
\newpage

%═══════════════════════════════════════════════════════════════
\section{Vue d'ensemble du projet}
%═══════════════════════════════════════════════════════════════

Le projet est organisé en deux grandes parties~:

\begin{center}
\begin{tikzpicture}[
  node distance=2cm,
  block/.style={draw=suva_blue, rounded corners, fill=blue!5,
    minimum width=4.5cm, minimum height=1cm,
    text width=4.2cm, align=center, font=\small},
  arrow/.style={->, thick, suva_blue}
]
  \node[block] (py) {\textbf{Python (off-chain)}\\\footnotesize
    \texttt{python-ref/}\\
    Génération de clés, Signature};
  \node[block, right=3cm of py] (sol) {\textbf{Solidity (on-chain)}\\\footnotesize
    \texttt{src/}\\
    Vérification on-chain};
  \draw[arrow] (py) -- node[above, font=\footnotesize] {(pk, sig)} (sol);
  \node[above=0.5cm of py, font=\scriptsize\color{gray}]
    {Machine utilisateur};
  \node[above=0.5cm of sol, font=\scriptsize\color{gray}]
    {Blockchain Ethereum};
\end{tikzpicture}
\end{center}

\bigskip
Structure du dossier principal \texttt{ETHDILITHIUM-main/}~:
\begin{lstlisting}[style=bash]
ETHDILITHIUM-main/
  makefile                  <- Commandes principales
  foundry.toml              <- Config Solidity (Foundry)
  python-ref/               <- Partie Python
    makefile
    requirements.txt
    Falcon_dilithium_py/    <- Code source Python
      Falcon_dilithium/     <- Algorithme Dilithium
      modules/              <- Arithmetique matricielle
      polynomials/          <- Anneaux polynomiaux
      Falcon_keccak_prng/   <- PRNG Keccak256
      shake/                <- Wrapper SHAKE
      utilities/            <- Utilitaires
      tests/                <- Tests unitaires
      benchmarks/           <- Benchmarks perf
  src/                      <- Contrats Solidity
  test/                     <- Tests Foundry (Solidity)
  lib/                      <- Dependances Foundry
\end{lstlisting}

%═══════════════════════════════════════════════════════════════
\section{Installation de la partie Python}
%═══════════════════════════════════════════════════════════════

\subsection{Prérequis}

\begin{itemize}
  \item Python 3.10 ou supérieur (\texttt{python3 --version})
  \item \texttt{pip} disponible
  \item \texttt{git} installé
\end{itemize}

\subsection{Étape 1 — Se placer dans le bon dossier}

\begin{cmdbox}{cd ETHDILITHIUM-main/python-ref}
\begin{lstlisting}[style=bash]
cd d:/Projects/suva/falcon_dilithium/ETHDILITHIUM-main/python-ref
\end{lstlisting}
\end{cmdbox}

\subsection{Étape 2 — Créer l'environnement virtuel}

\begin{cmdbox}{Créer le venv}
\begin{lstlisting}[style=bash]
python3 -m venv myenv
\end{lstlisting}
\end{cmdbox}

\subsection{Étape 3 — Activer l'environnement}

\begin{tcolorbox}[colback=gray!5, colframe=gray!40, title=\textbf{Windows}]
\begin{lstlisting}[style=bash]
myenv\Scripts\activate
\end{lstlisting}
\end{tcolorbox}

\begin{tcolorbox}[colback=gray!5, colframe=gray!40, title=\textbf{Linux / macOS / Git Bash}]
\begin{lstlisting}[style=bash]
source myenv/bin/activate
\end{lstlisting}
\end{tcolorbox}

Ton prompt doit afficher \texttt{(myenv)} au début.

\subsection{Étape 4 — Installer les dépendances}

\begin{cmdbox}{Installer les packages}
\begin{lstlisting}[style=bash]
pip install -r requirements.txt
pip install -e "git+https://github.com/ZKNoxHQ/NTT.git@main#egg=polyntt&subdirectory=assets/pythonref/"
\end{lstlisting}
\end{cmdbox}

\begin{infobox}
Le package \texttt{polyntt} est le moteur NTT développé par ZKNox.
Il est installé directement depuis GitHub. Tu as besoin d'une connexion internet.
\end{infobox}

\subsection{Étape 5 — Vérifier l'installation}

\begin{cmdbox}{Test rapide d'import}
\begin{lstlisting}[style=bash]
python3 -c "import polyntt; print('polyntt OK')"
python3 -c "from Crypto.Cipher import AES; print('pycryptodome OK')"
\end{lstlisting}
\end{cmdbox}

\begin{successbox}
\begin{lstlisting}[style=output]
polyntt OK
pycryptodome OK
\end{lstlisting}
\end{successbox}

%═══════════════════════════════════════════════════════════════
\section{Lancer les tests Python}
%═══════════════════════════════════════════════════════════════

\subsection{Important --- le nom du package}

\begin{warningbox}
Le dossier s'appelle \texttt{Falcon\_dilithium\_py/} mais les tests importent
depuis \texttt{dilithium\_py}. Pour résoudre ça, lance les tests
\textbf{depuis le dossier} \texttt{python-ref/} en pointant vers
\texttt{Falcon\_dilithium\_py} directement.
\end{warningbox}

\subsection{Méthode recommandée --- unittest discover}

Se placer dans \texttt{python-ref/Falcon\_dilithium\_py/} puis~:

\begin{cmdbox}{Lancer TOUS les tests}
\begin{lstlisting}[style=bash]
cd Falcon_dilithium_py
python3 -m pytest tests/ -v
\end{lstlisting}
\end{cmdbox}

Ou avec unittest~:
\begin{cmdbox}{Avec unittest}
\begin{lstlisting}[style=bash]
python3 -m unittest discover -s tests -v
\end{lstlisting}
\end{cmdbox}

\subsection{Lancer un test spécifique}

\begin{cmdbox}{Un seul fichier de test}
\begin{lstlisting}[style=bash]
python3 -m pytest tests/test_dilithium.py -v
python3 -m pytest tests/test_utils.py -v
python3 -m pytest tests/test_polynomial_generic.py -v
\end{lstlisting}
\end{cmdbox}

\subsection{Description de chaque fichier de test}

\begin{center}
\begin{tabular}{lp{9cm}}
\toprule
\textbf{Fichier} & \textbf{Ce qu'il teste} \\
\midrule
\texttt{test\_dilithium.py} &
  Le flux complet : keygen → sign → verify.
  Teste aussi que la vérification échoue avec une mauvaise clé/message.
  Inclut les KAT (Known Answer Tests) NIST. \\
\addlinespace
\texttt{test\_drbg.py} &
  Le générateur pseudo-aléatoire AES256-CTR.
  Vérifie la déterminisme : même seed → mêmes bytes. \\
\addlinespace
\texttt{test\_polynomial\_generic.py} &
  Les opérations sur les polynômes dans $R_q$.
  Addition, soustraction, multiplication via NTT. \\
\addlinespace
\texttt{test\_module\_generic.py} &
  La multiplication matrice-vecteur $A \cdot s_1$.
  Operations sur les modules (vecteurs de polynômes). \\
\addlinespace
\texttt{test\_shake.py} &
  Le wrapper SHAKE-128 et SHAKE-256. \\
\addlinespace
\texttt{test\_utils.py} &
  Les fonctions utilitaires : \texttt{decompose},
  \texttt{high\_bits}, \texttt{make\_hint}, \texttt{use\_hint}. \\
\bottomrule
\end{tabular}
\end{center}

\subsection{Sortie attendue d'un test réussi}

\begin{successbox}
\begin{lstlisting}[style=output]
test_dilithium2 ... ok
test_dilithium2_keccak_prng ... ok
test_ethdilithium2_keccak_prng ... ok
test_ethdilithium2_shake ... ok

----------------------------------------------------------------------
Ran 4 tests in 12.453s

OK
\end{lstlisting}
\end{successbox}

%═══════════════════════════════════════════════════════════════
\section{Lancer les exemples}
%═══════════════════════════════════════════════════════════════

\subsection{Exemple de base}

\begin{cmdbox}{Exemple standard Dilithium}
\begin{lstlisting}[style=bash]
cd python-ref/Falcon_dilithium_py
python3 example.py
\end{lstlisting}
\end{cmdbox}

Ce script fait~:
\begin{lstlisting}[style=python]
from Falcon_dilithium import ETHDilithium2, Dilithium2

msg = b"We are ZKNox."

# Dilithium standard
pk, sk = Dilithium2.keygen()
sig = Dilithium2.sign(sk, msg)
assert Dilithium2.verify(pk, msg, sig)
print("Dilithium2 OK")

# ETHDilithium (optimise Keccak)
pk, sk = ETHDilithium2.keygen()
sig = ETHDilithium2.sign(sk, msg)
assert ETHDilithium2.verify(pk, msg, sig)
print("ETHDilithium2 OK")
\end{lstlisting}

\subsection{Exemple ZK (BabyBear / KoalaBear)}

\begin{cmdbox}{Exemple ZK}
\begin{lstlisting}[style=bash]
python3 example_zk.py
\end{lstlisting}
\end{cmdbox}

Ce script teste les variantes ZK-friendly avec des corps alternatifs.

%═══════════════════════════════════════════════════════════════
\section{Test manuel pas-à-pas en Python interactif}
%═══════════════════════════════════════════════════════════════

Tu peux tester chaque composant séparément depuis le REPL Python.
C'est la meilleure façon de comprendre ce qui se passe.

\subsection{Test de l'arithmétique polynomiale}

\begin{lstlisting}[style=python]
# Depuis python-ref/Falcon_dilithium_py/
import sys
sys.path.insert(0, '.')

from polynomials.polynomials import PolynomialRingDilithium

# Creer l'anneau R_q = Z_q[x]/(x^256 + 1)
R = PolynomialRingDilithium()

# Creer deux polynomes aleatoires
p1 = R.random_element()  # coefficients aleatoires dans Z_q
p2 = R.random_element()

# Tester les operations
p3 = p1 + p2      # addition
p4 = p1 * p2      # multiplication (via NTT)
p5 = p1 - p2      # soustraction

print(f"Coefficients de p1[0:5] = {p1.coeffs[:5]}")
print(f"Coefficients de p3[0:5] = {p3.coeffs[:5]}")
print(f"Multiplication OK ? {p4.coeffs[0] != 0}")
\end{lstlisting}

\subsection{Test du keygen complet}

\begin{lstlisting}[style=python]
from Falcon_dilithium.Falcon_dilithium import Dilithium

# Parametres Dilithium2
dil = Dilithium(...)  # voir Falcon_default_parameters.py

# Generer une paire de cles
pk, sk = dil.keygen()
print(f"Taille cle publique : {len(pk)} bytes")
# Attendu : 1312 bytes pour Dilithium2

print(f"Taille cle secrete : {len(sk)} bytes")
# Attendu : 2528 bytes pour Dilithium2

# Signer un message
msg = b"Hello SuVa"
sig = dil.sign(sk, msg)
print(f"Taille signature : {len(sig)} bytes")
# Attendu : 2420 bytes pour Dilithium2

# Verifier
ok = dil.verify(pk, msg, sig)
print(f"Verification : {ok}")  # True
\end{lstlisting}

\subsection{Test des utilitaires}

\begin{lstlisting}[style=python]
from utilities.utils import decompose, high_bits, low_bits, make_hint, use_hint

q = 8380417
alpha = 2 * ((q - 1) // 88)  # 2 * gamma_2

# Decomposer un element
r = 1234567
r1, r0 = decompose(r, alpha)
print(f"r  = {r}")
print(f"r1 = {r1}  (bits hauts)")
print(f"r0 = {r0}  (bits bas)")
print(f"Verification : r1*alpha + r0 = {r1*alpha + r0}")

# Le hint
h = make_hint(r0, r1, alpha)
r1_recovered = use_hint(h, r, alpha)
print(f"r1 original   = {r1}")
print(f"r1 recupere   = {r1_recovered}")
print(f"Hint fonctionne : {r1 == r1_recovered}")
\end{lstlisting}

%═══════════════════════════════════════════════════════════════
\section{Benchmarks}
%═══════════════════════════════════════════════════════════════

\begin{cmdbox}{Lancer les benchmarks}
\begin{lstlisting}[style=bash]
cd python-ref/Falcon_dilithium_py
python3 -m benchmarks.benchmark_dilithium
\end{lstlisting}
\end{cmdbox}

\begin{successbox}
Résultats typiques sur une machine moderne (1000 itérations)~:
\begin{lstlisting}[style=output]
Dilithium2
  keygen  : median =  6.2 ms
  sign    : median = 27.4 ms  (moyenne 27.8 ms)
  verify  : median =  7.1 ms

Dilithium3
  keygen  : median =  9.5 ms
  sign    : median = 46.1 ms
  verify  : median = 11.3 ms
\end{lstlisting}
\end{successbox}

%═══════════════════════════════════════════════════════════════
\section{Installation et tests Solidity (Foundry)}
%═══════════════════════════════════════════════════════════════

\subsection{Installer Foundry}

\begin{cmdbox}{Installation de Foundry}
\begin{lstlisting}[style=bash]
# Installer foundryup (installer de Foundry)
curl -L https://foundry.paradigm.xyz | bash
# Puis dans un nouveau terminal :
foundryup
\end{lstlisting}
\end{cmdbox}

\begin{infobox}
Foundry installe~: \texttt{forge} (compilateur/testeur),
\texttt{cast} (interactions avec la blockchain),
\texttt{anvil} (nœud local).
Vérifie avec \texttt{forge --version}.
\end{infobox}

\subsection{Générer les vecteurs de test}

Les tests Solidity nécessitent d'abord de générer des vecteurs de test
depuis Python~:

\begin{cmdbox}{Générer les vecteurs de test}
\begin{lstlisting}[style=bash]
cd ETHDILITHIUM-main/python-ref
# Generer les vecteurs ETHDilithium
python3 -m Falcon_dilithium_py.generate_ethdilithium_test_vectors
# Generer les vecteurs Dilithium standard
python3 -m Falcon_dilithium_py.generate_dilithium_test_vectors
\end{lstlisting}
\end{cmdbox}

Cela écrit un fichier \texttt{test/ZKNOX\_dilithium.t.sol} avec les
données de test en dur pour Foundry.

\subsection{Lancer les tests Solidity}

\begin{cmdbox}{Tests Solidity rapides (profil lite)}
\begin{lstlisting}[style=bash]
cd ETHDILITHIUM-main
FOUNDRY_PROFILE=lite forge test -vv
\end{lstlisting}
\end{cmdbox}

\begin{cmdbox}{Tests Solidity complets (plus lents)}
\begin{lstlisting}[style=bash]
forge test -vv
\end{lstlisting}
\end{cmdbox}

\begin{infobox}
\texttt{-vv} = verbose level 2 (affiche les logs des tests).\\
\texttt{-vvv} = encore plus verbeux (affiche les traces d'exécution).
\end{infobox}

\subsection{Vérification gas}

Pour voir le coût en gas de chaque opération~:

\begin{cmdbox}{Rapport gas}
\begin{lstlisting}[style=bash]
forge test --gas-report
\end{lstlisting}
\end{cmdbox}

%═══════════════════════════════════════════════════════════════
\section{Erreurs fréquentes et solutions}
%═══════════════════════════════════════════════════════════════

\begin{errorbox}
\textbf{Erreur~:} \texttt{ModuleNotFoundError: No module named 'polyntt'}

\textbf{Solution~:} Le package NTT n'est pas installé.
\begin{lstlisting}[style=bash]
pip install -e "git+https://github.com/ZKNoxHQ/NTT.git@main#egg=polyntt&subdirectory=assets/pythonref/"
\end{lstlisting}
\end{errorbox}

\begin{errorbox}
\textbf{Erreur~:} \texttt{ModuleNotFoundError: No module named 'Crypto'}

\textbf{Solution~:} pycryptodome mal installé.
\begin{lstlisting}[style=bash]
pip uninstall pycryptodome pycrypto
pip install pycryptodome==3.14.1
\end{lstlisting}
\end{errorbox}

\begin{errorbox}
\textbf{Erreur~:} \texttt{ImportError: attempted relative import beyond top-level package}

\textbf{Solution~:} Lance Python depuis le bon répertoire.
Être dans \texttt{python-ref/} et utiliser \texttt{python3 -m Falcon\_dilithium\_py.tests.test\_dilithium}.
\end{errorbox}

\begin{errorbox}
\textbf{Erreur~:} \texttt{forge: command not found}

\textbf{Solution~:} Foundry n'est pas dans le PATH.
\begin{lstlisting}[style=bash]
# Ajouter ~/.foundry/bin au PATH (Linux/macOS)
export PATH="$HOME/.foundry/bin:$PATH"
# Sur Windows, redemarrer le terminal apres foundryup
\end{lstlisting}
\end{errorbox}

%═══════════════════════════════════════════════════════════════
\section{Récapitulatif des commandes clés}
%═══════════════════════════════════════════════════════════════

\begin{center}
\begin{tabular}{ll}
\toprule
\textbf{Action} & \textbf{Commande} \\
\midrule
Installer Python deps & \texttt{make install} (dans \texttt{python-ref/}) \\
Lancer tests Python & \texttt{make test} (dans \texttt{python-ref/}) \\
Lancer exemple & \texttt{make example} \\
Générer test vectors & \texttt{make generate\_test\_vectors} \\
Installer Foundry & \texttt{foundryup} \\
Tests Solidity rapides & \texttt{FOUNDRY\_PROFILE=lite forge test -vv} \\
Tests Solidity complets & \texttt{forge test -vv} \\
Rapport gas & \texttt{forge test --gas-report} \\
Compiler Solidity & \texttt{forge build} \\
\bottomrule
\end{tabular}
\end{center}

\end{document}

\documentclass[12pt, a4paper]{article}
\usepackage[utf8]{inputenc}
\usepackage[T1]{fontenc}
\usepackage[french]{babel}
\usepackage[margin=2.5cm]{geometry}
\usepackage{amsmath, amssymb, amsthm, mathtools}
\usepackage{listings}
\usepackage{xcolor}
\usepackage{tcolorbox}
\tcbuselibrary{theorems, skins, breakable}
\usepackage{booktabs}
\usepackage{array}
\usepackage{fancyhdr}
\usepackage{hyperref}
\usepackage{tikz}
\usetikzlibrary{shapes, positioning}

% ── Couleurs ──────────────────────────────────────────────────────────
\definecolor{suva_blue}{RGB}{30, 100, 180}
\definecolor{code_bg}{RGB}{245, 247, 250}
\definecolor{code_kw}{RGB}{0, 80, 200}
\definecolor{code_str}{RGB}{180, 40, 40}
\definecolor{code_cmt}{RGB}{80, 130, 80}
\definecolor{warn_bg}{RGB}{255, 248, 220}
\definecolor{success_bg}{RGB}{230, 255, 235}
\definecolor{dark_bg}{RGB}{25, 30, 50}

\lstdefinestyle{python}{
  language=Python,
  basicstyle=\ttfamily\small,
  keywordstyle=\color{code_kw}\bfseries,
  stringstyle=\color{code_str},
  commentstyle=\color{code_cmt}\itshape,
  backgroundcolor=\color{code_bg},
  frame=single, rulecolor=\color{gray!40},
  numbers=left, numberstyle=\tiny\color{gray},
  breaklines=true, tabsize=4, showstringspaces=false,
}
\lstdefinestyle{output}{
  basicstyle=\ttfamily\small,
  backgroundcolor=\color{gray!10},
  frame=single, rulecolor=\color{gray!50},
  breaklines=true,
}

\newtcolorbox{defbox}[2][]{
  colback=success_bg, colframe=green!50!black,
  fonttitle=\bfseries, title={Définition : #2}, #1
}
\newtcolorbox{theorembox}[2][]{
  colback=blue!5, colframe=suva_blue,
  fonttitle=\bfseries, title={Théorème : #2}, #1
}
\newtcolorbox{exercice}[2][]{
  colback=white, colframe=suva_blue!60,
  fonttitle=\bfseries, title={Exercice #2}, #1, breakable
}
\newtcolorbox{solution}[1][]{
  colback=gray!5, colframe=gray!40,
  fonttitle=\bfseries\itshape, title={Solution}, #1, breakable
}
\newtcolorbox{importantbox}[1][]{
  colback=suva_blue!10, colframe=suva_blue,
  fonttitle=\bfseries, title={Connexion avec le projet}, #1
}
\newtcolorbox{warningbox}[1][]{
  colback=warn_bg, colframe=orange!80!black,
  fonttitle=\bfseries, title={Attention}, #1
}

\newcommand{\Zq}{\mathbb{Z}_q}
\newcommand{\ZZ}{\mathbb{Z}}
\newcommand{\pmod}[1]{\ (\mathrm{mod}\ #1)}

\newcommand{\auteur}{Mohamed Lamine MBENGUE}
\newcommand{\specialites}{Sécurité · DevOps · Dev Full Stack Web \& Mobile}
\newcommand{\tel}{+221 78 178 44 36}
\newcommand{\email}{mbenguemohamedlamine2022@gmail.com}
\newcommand{\portfolio}{portfolio-alamine.web.app}

\pagestyle{fancy}
\fancyhf{}
\fancyhead[L]{\color{suva_blue}\textbf{SuVa --- Chapitre 1}}
\fancyhead[R]{\color{gray}Arithmétique Modulaire}
\fancyfoot[L]{\footnotesize\color{gray}\auteur}
\fancyfoot[C]{\thepage}
\fancyfoot[R]{\footnotesize\color{gray}\portfolio}
\renewcommand{\headrulewidth}{0.4pt}
\renewcommand{\footrulewidth}{0.3pt}

% ══════════════════════════════════════════════════════════════════════
\begin{document}

\begin{center}
  {\color{suva_blue}\rule{\linewidth}{2pt}}\\[0.5cm]
  {\Huge\bfseries\color{suva_blue} Chapitre 1}\\[0.3cm]
  {\large\bfseries Arithmétique Modulaire et Corps Finis}\\[0.2cm]
  {\small\color{gray} Fondement mathématique de toute la cryptographie du projet}\\[0.3cm]
  {\color{suva_blue}\rule{\linewidth}{2pt}}
\end{center}

\begin{tcolorbox}[colback=suva_blue!5, colframe=suva_blue!40,
  left=6pt, right=6pt, top=4pt, bottom=4pt]
\small
\begin{tabular}{ll}
  \textbf{Auteur} & \auteur \\
  \textbf{Spécialités} & \specialites \\
  \textbf{Tél} & \tel \quad \textbf{Email} : \email \\
  \textbf{Portfolio} & \href{https://\portfolio}{\texttt{\portfolio}} \\
\end{tabular}
\end{tcolorbox}

\tableofcontents
\newpage

%══════════════════════════════════════════════════════════════════════
\section{Arithmétique modulaire --- rappels}
%══════════════════════════════════════════════════════════════════════

\subsection{L'opération modulo}

\begin{defbox}{Modulo}
Pour deux entiers $a \in \ZZ$ et $m \in \ZZ_{>0}$, on dit que
$a \equiv r \pmod{m}$ si $m$ divise $a - r$.
L'entier $r$ est le \textbf{reste} de la division euclidienne de $a$ par $m$.

On note aussi $a \bmod m = r$ avec $r \in \{0, 1, \ldots, m-1\}$.
\end{defbox}

\textbf{Exemples~:}
\begin{align*}
  17 \bmod 5 &= 2  &\text{car } 17 = 3 \times 5 + 2\\
  -3 \bmod 7 &= 4  &\text{car } -3 = (-1) \times 7 + 4\\
  100 \bmod 13 &= 9 &\text{car } 100 = 7 \times 13 + 9
\end{align*}

\begin{warningbox}
En Python, \texttt{-3 \% 7 = 4} (correct). En C, le résultat peut être négatif.
Dans le projet, tous les calculs doivent donner un résultat dans $[0, q-1]$.
\end{warningbox}

\subsection{L'anneau $\mathbb{Z}_q$}

\begin{defbox}{L'anneau $\mathbb{Z}_q$}
$\mathbb{Z}_q = \{0, 1, 2, \ldots, q-1\}$ muni de l'addition et de la
multiplication modulo $q$. C'est un \textbf{corps} (corps fini) si et
seulement si $q$ est \textbf{premier}.
\end{defbox}

\begin{importantbox}
Dans le projet SuVa / Dilithium, le module est~:
$$q = 8\,380\,417$$
Ce nombre est premier (vérifie-le !). Tous les coefficients des polynômes
vivent dans $\mathbb{Z}_q = \{0, 1, \ldots, 8\,380\,416\}$.
\end{importantbox}

\subsection{Propriétés de l'addition et de la multiplication modulo $q$}

Les règles habituelles s'appliquent~:
\begin{align*}
  (a + b) \bmod q &\equiv ((a \bmod q) + (b \bmod q)) \bmod q\\
  (a \cdot b) \bmod q &\equiv ((a \bmod q) \cdot (b \bmod q)) \bmod q
\end{align*}

\textbf{Exemple avec $q = 8\,380\,417$~:}
\begin{lstlisting}[style=python]
q = 8_380_417

a = 5_000_000
b = 4_000_000

addition = (a + b) % q
print(f"({a} + {b}) mod q = {addition}")
# (5000000 + 4000000) mod 8380417 = 619583

multiplication = (a * b) % q
print(f"({a} * {b}) mod q = {multiplication}")
# Resultat dans [0, 8380416]
\end{lstlisting}

%══════════════════════════════════════════════════════════════════════
\section{L'inverse modulaire}
%══════════════════════════════════════════════════════════════════════

\subsection{Définition}

\begin{defbox}{Inverse modulaire}
L'inverse de $a$ modulo $q$ est l'entier $a^{-1} \in \{1, \ldots, q-1\}$
tel que~:
$$a \cdot a^{-1} \equiv 1 \pmod{q}$$
L'inverse existe si et seulement si $\gcd(a, q) = 1$.
Si $q$ est premier, tout $a \neq 0$ est inversible.
\end{defbox}

\subsection{Algorithme d'Euclide étendu}

\begin{theorembox}{Théorème de Bézout}
Pour tout $a, b$ tels que $\gcd(a, b) = d$, il existe $u, v \in \ZZ$
tels que~:
$$au + bv = d$$
Si $d = 1$, alors $u = a^{-1} \pmod{b}$.
\end{theorembox}

\begin{lstlisting}[style=python]
def xgcd(a, b):
    """
    Algorithme d'Euclide etendu.
    Retourne (gcd, u, v) tels que a*u + b*v = gcd.
    """
    if b == 0:
        return a, 1, 0
    g, u, v = xgcd(b, a % b)
    return g, v, u - (a // b) * v

def mod_inverse(a, q):
    """Inverse de a modulo q (q premier)."""
    g, u, v = xgcd(a % q, q)
    assert g == 1, f"{a} n'est pas inversible modulo {q}"
    return u % q

# Test
q = 8_380_417
a = 42
inv = mod_inverse(a, q)
print(f"42^(-1) mod q = {inv}")
print(f"Verification : 42 * {inv} mod q = {(42 * inv) % q}")  # Doit etre 1
\end{lstlisting}

\subsection{Méthode rapide : petit théorème de Fermat}

\begin{theorembox}{Petit théorème de Fermat}
Si $q$ est premier et $\gcd(a, q) = 1$, alors~:
$$a^{q-1} \equiv 1 \pmod{q}$$
Donc~: $a^{-1} \equiv a^{q-2} \pmod{q}$
\end{theorembox}

\begin{lstlisting}[style=python]
def mod_inverse_fermat(a, q):
    """Inverse via le petit theoreme de Fermat (q premier)."""
    return pow(a, q - 2, q)  # pow(a, e, m) = a^e mod m (Python integre)

q = 8_380_417
a = 42
inv = mod_inverse_fermat(a, q)
print(f"42^(-1) mod q = {inv}")
print(f"Verification  : {(42 * inv) % q}")  # 1
\end{lstlisting}

\begin{importantbox}
La fonction \texttt{pow(a, q-2, q)} de Python est \textbf{très efficace}~:
elle utilise l'exponentiation rapide en $O(\log q)$ multiplications.
C'est ce qui est utilisé dans le projet pour les inversions modulaires.
\end{importantbox}

%══════════════════════════════════════════════════════════════════════
\section{Réduction symétrique (centred reduction)}
%══════════════════════════════════════════════════════════════════════

\subsection{Motivation}

Dans Dilithium, les coefficients ne vivent pas toujours dans $[0, q-1]$.
Parfois on préfère les représenter dans $[-(q-1)/2, (q-1)/2]$.

\begin{defbox}{Réduction symétrique}
La réduction symétrique de $a$ modulo $q$ est~:
$$a_{\mathrm{sym}} = \begin{cases}
  a \bmod q & \text{si } (a \bmod q) \leq (q-1)/2\\
  (a \bmod q) - q & \text{sinon}
\end{cases}$$
Résultat dans $[-(q-1)/2, (q-1)/2]$.
\end{defbox}

\begin{lstlisting}[style=python]
# Fichier : utilities/utils.py dans le projet
def reduce_mod_pm(a, q):
    """Reduction symmetrique de a modulo q."""
    r = a % q
    if r > (q - 1) // 2:
        r -= q
    return r

q = 8_380_417
# Exemples
print(reduce_mod_pm(8_380_416, q))  # = -1 (pas q-1)
print(reduce_mod_pm(1, q))          # =  1
print(reduce_mod_pm(q // 2 + 1, q)) # = -(q//2)
\end{lstlisting}

\begin{importantbox}
Cette fonction apparaît dans \texttt{utilities/utils.py}.
Elle est utilisée dans la vérification de normes ($\|z\|_\infty < \gamma_1 - \beta$).
\end{importantbox}

%══════════════════════════════════════════════════════════════════════
\section{Décomposition power-of-2 (HighBits / LowBits)}
%══════════════════════════════════════════════════════════════════════

C'est un mécanisme crucial dans Dilithium pour compresser la clé publique.

\subsection{Principe}

\begin{defbox}{Décomposition modulaire}
Pour $r \in \ZZ_q$ et $\alpha$ pair divisant $q - 1$, on écrit~:
$$r = r_1 \cdot \alpha + r_0$$
avec $r_0 \in (- \alpha/2,\ \alpha/2]$ (partie basse, \emph{LowBits})
et $r_1 = (r - r_0) / \alpha$ (partie haute, \emph{HighBits}).
\end{defbox}

\textbf{Pourquoi~?} Dans Dilithium, $t = A s_1 + s_2$ est un vecteur de polynômes.
On décompose $t = t_1 \cdot 2^d + t_0$. La clé publique ne publie que $t_1$
(bits hauts), économisant $d = 13$ bits par coefficient. $t_0$ reste secret.

\begin{lstlisting}[style=python]
# Extrait de utilities/utils.py
def decompose(r, alpha, q=8_380_417):
    """
    Decomposes r = r1 * alpha + r0
    r0 in (-alpha/2, alpha/2]
    """
    r = r % q
    r0 = r % alpha
    if r0 > alpha // 2:
        r0 -= alpha
    r1 = (r - r0) // alpha
    if r1 == (q - 1) // alpha + 1:
        r1 = 0
        r0 -= 1
    return r1, r0

def high_bits(r, alpha, q=8_380_417):
    r1, _ = decompose(r, alpha, q)
    return r1

def low_bits(r, alpha, q=8_380_417):
    _, r0 = decompose(r, alpha, q)
    return r0

# Test
q = 8_380_417
d = 13
alpha = 1 << d  # 2^13 = 8192

r = 100_000
r1, r0 = decompose(r, alpha, q)
print(f"r  = {r}")
print(f"r1 = {r1}  (partie haute, dans la cle publique)")
print(f"r0 = {r0}  (partie basse, secrete)")
print(f"Verification : r1 * {alpha} + r0 = {r1 * alpha + r0}")
# Doit etre congruent a r mod q
\end{lstlisting}

%══════════════════════════════════════════════════════════════════════
\section{Le vecteur hint --- make\_hint et use\_hint}
%══════════════════════════════════════════════════════════════════════

\subsection{Motivation}

Lors de la vérification, on recalcule $w' = Az - ct_1 \cdot 2^d$ et on
veut retrouver $w_1 = \mathrm{HighBits}(Ay)$ d'origine.
Le problème~: $Az - ct_1 \cdot 2^d \approx Ay$ mais pas exactement.
Le \textbf{hint} $h$ encode la correction nécessaire de manière compacte.

\begin{defbox}{Hint}
Pour $r \in \ZZ_q$ et $r + z \in \ZZ_q$ avec $z$ petit~:
$$h = \mathrm{MakeHint}(z, r, \alpha) \in \{0, 1\}$$
$$r_1 = \mathrm{UseHint}(h, r + z, \alpha)$$
de sorte que $r_1 = \mathrm{HighBits}(r, \alpha)$.
\end{defbox}

\begin{lstlisting}[style=python]
# Extrait de utilities/utils.py
def make_hint(z, r, alpha, q=8_380_417):
    """
    Cree un hint indiquant si HighBits(r, alpha) != HighBits(r+z, alpha).
    Retourne 0 ou 1.
    """
    r1 = high_bits(r, alpha, q)
    v1 = high_bits(r + z, alpha, q)
    return int(r1 != v1)

def use_hint(h, r, alpha, q=8_380_417):
    """
    Utilise le hint pour recuperer HighBits(r - z, alpha)
    sans connaitre z.
    """
    m = (q - 1) // alpha
    r1, r0 = decompose(r, alpha, q)
    if h == 1 and r0 > 0:
        return (r1 + 1) % m
    if h == 1 and r0 <= 0:
        return (r1 - 1) % m
    return r1

# Demonstration
q = 8_380_417
alpha = 2 * ((q - 1) // 88)  # = 2 * gamma_2 dans Dilithium

r = 500_000   # valeur originale
z = 90_000    # petite perturbation (typiquement cs2 + ct0)

r1_original = high_bits(r, alpha, q)
h = make_hint(z, r, alpha, q)
r1_recovered = use_hint(h, r + z, alpha, q)

print(f"r1 original  = {r1_original}")
print(f"hint h       = {h}")
print(f"r1 recupere  = {r1_recovered}")
print(f"Correct ?    = {r1_original == r1_recovered}")
\end{lstlisting}

%══════════════════════════════════════════════════════════════════════
\section{Vérification de la norme infinie}
%══════════════════════════════════════════════════════════════════════

\begin{defbox}{Norme infinie}
Pour un polynôme $f = \sum_{i=0}^{n-1} f_i \cdot x^i \in R_q$, la norme
infinie (après réduction symétrique) est~:
$$\|f\|_\infty = \max_{i=0}^{n-1} |f_i|$$
\end{defbox}

\begin{importantbox}
Dans Dilithium, les rejections sampling vérifient~:
\begin{itemize}
  \item $\|z\|_\infty < \gamma_1 - \beta$ \hfill ($\gamma_1 = 2^{17}$, $\beta = 78$)
  \item $\|w_0 - cs_2\|_\infty < \gamma_2 - \beta$ \hfill ($\gamma_2 = (q-1)/88$)
  \item $\|ct_0\|_\infty < \gamma_2$
\end{itemize}
Si une de ces conditions échoue, le signe recommence.
\end{importantbox}

\begin{lstlisting}[style=python]
def check_norm_bound(poly_coeffs, bound, q=8_380_417):
    """
    Verifie si |coefficient| >= bound pour au moins un coefficient.
    Retourne True si la norme depasse la borne (==> rejeter).
    """
    for coeff in poly_coeffs:
        # Reduction symmetrique
        c = coeff % q
        if c > (q - 1) // 2:
            c -= q
        if abs(c) >= bound:
            return True
    return False

# Test
import random
q = 8_380_417
gamma_1 = 1 << 17    # 131072
beta = 78
bound = gamma_1 - beta  # 130994

# Un vecteur z valide (petits coefficients)
z_valid = [random.randint(-bound + 1, bound - 1) for _ in range(256)]
print(f"z_valid depasse borne ? {check_norm_bound(z_valid, bound)}")  # False

# Un vecteur z invalide (trop grand)
z_bad = z_valid.copy()
z_bad[42] = gamma_1  # coefficient trop grand
print(f"z_bad depasse borne  ? {check_norm_bound(z_bad, bound)}")   # True
\end{lstlisting}

%══════════════════════════════════════════════════════════════════════
\section{Exercices}
%══════════════════════════════════════════════════════════════════════

\begin{exercice}{1 --- Calculs modulo q}
Soit $q = 8\,380\,417$.

\begin{enumerate}[label=\alph*)]
  \item Calcule $(7\,000\,000 + 2\,000\,000) \bmod q$.
  \item Calcule $(-1) \bmod q$.
  \item Calcule $q^2 \bmod q$.
  \item Trouve $a^{-1} \bmod q$ pour $a = 256$.
  \item Vérifie que $256 \times a^{-1} \equiv 1 \pmod{q}$.
\end{enumerate}
\end{exercice}

\begin{solution}
\begin{lstlisting}[style=python]
q = 8_380_417

# a)
print((7_000_000 + 2_000_000) % q)   # 619583

# b)
print((-1) % q)                        # 8380416

# c)
print((q**2) % q)                      # 0  (q^2 = q * q)

# d)
inv_256 = pow(256, q - 2, q)
print(f"256^(-1) mod q = {inv_256}")

# e)
print((256 * inv_256) % q)             # 1
\end{lstlisting}
\end{solution}

\begin{exercice}{2 --- Reduction symétrique}
\begin{enumerate}[label=\alph*)]
  \item Calcule la réduction symétrique de $8\,380\,416$ modulo $q$.
  \item Calcule la réduction symétrique de $4\,000\,000$ modulo $q$.
  \item Quelle est la plage de la réduction symétrique pour $q = 8\,380\,417$ ?
\end{enumerate}
\end{exercice}

\begin{solution}
\begin{lstlisting}[style=python]
q = 8_380_417

def sym(a, q):
    r = a % q
    return r - q if r > (q-1)//2 else r

print(sym(8_380_416, q))  # -1
print(sym(4_000_000, q))  # 4000000 (deja dans la plage)
# Plage : [-(q-1)//2, (q-1)//2] = [-4190208, 4190208]
\end{lstlisting}
\end{solution}

\begin{exercice}{3 --- Décomposition et hint (le plus important)}
Soit $q = 8\,380\,417$, $\alpha = 2\gamma_2$ avec $\gamma_2 = (q-1)/88$.

\begin{enumerate}[label=\alph*)]
  \item Calcule $\alpha$.
  \item Décompose $r = 3\,000\,000$ : calcule $r_1, r_0$.
  \item Vérifie que $r_1 \times \alpha + r_0 \equiv r \pmod{q}$.
  \item Soit $z = 50\,000$. Calcule $h = \mathrm{MakeHint}(z, r, \alpha)$.
  \item Vérifie que $\mathrm{UseHint}(h, r+z, \alpha) = \mathrm{HighBits}(r, \alpha)$.
\end{enumerate}
\end{exercice}

\begin{solution}
\begin{lstlisting}[style=python]
q = 8_380_417
gamma_2 = (q - 1) // 88
alpha = 2 * gamma_2
print(f"alpha = {alpha}")  # 190464

# b) Decomposer r = 3_000_000
r = 3_000_000
r0 = r % alpha
if r0 > alpha // 2:
    r0 -= alpha
r1 = (r - r0) // alpha
print(f"r1 = {r1}, r0 = {r0}")

# c) Verification
print(f"r1*alpha + r0 = {r1 * alpha + r0}")  # == r

# d) Hint
z = 50_000
h = make_hint(z, r, alpha, q)
print(f"hint h = {h}")

# e) UseHint
r1_rec = use_hint(h, r + z, alpha, q)
r1_orig = high_bits(r, alpha, q)
print(f"r1_orig = {r1_orig}, r1_rec = {r1_rec}")
print(f"Correct : {r1_orig == r1_rec}")
\end{lstlisting}
\end{solution}

\begin{exercice}{4 --- Exploration du code source}
Ouvre le fichier \texttt{utilities/utils.py} du projet.
\begin{enumerate}[label=\alph*)]
  \item Identifie la différence entre \texttt{make\_hint} et \texttt{make\_hint\_optimised}.
        (Indice~: Dilithium 5.1 vs spec NIST)
  \item Quel est l'avantage de \texttt{make\_hint\_optimised} ?
  \item À quel endroit dans \texttt{Falcon\_dilithium.py} est appelé \texttt{make\_hint\_optimised} ?
\end{enumerate}
\end{exercice}

\end{document}

\documentclass[12pt, a4paper]{article}
\usepackage[utf8]{inputenc}
\usepackage[T1]{fontenc}
\usepackage[french]{babel}
\usepackage[margin=2.5cm]{geometry}
\usepackage{amsmath, amssymb, amsthm, mathtools}
\usepackage{listings}
\usepackage{xcolor}
\usepackage{tcolorbox}
\tcbuselibrary{theorems, skins, breakable}
\usepackage{booktabs}
\usepackage{fancyhdr}
\usepackage{hyperref}
\usepackage{tikz}

\definecolor{suva_blue}{RGB}{30, 100, 180}
\definecolor{code_bg}{RGB}{245, 247, 250}
\definecolor{code_kw}{RGB}{0, 80, 200}
\definecolor{code_str}{RGB}{180, 40, 40}
\definecolor{code_cmt}{RGB}{80, 130, 80}
\definecolor{warn_bg}{RGB}{255, 248, 220}
\definecolor{success_bg}{RGB}{230, 255, 235}

\lstdefinestyle{python}{
  language=Python,
  basicstyle=\ttfamily\small,
  keywordstyle=\color{code_kw}\bfseries,
  stringstyle=\color{code_str},
  commentstyle=\color{code_cmt}\itshape,
  backgroundcolor=\color{code_bg},
  frame=single, rulecolor=\color{gray!40},
  numbers=left, numberstyle=\tiny\color{gray},
  breaklines=true, tabsize=4, showstringspaces=false,
}
\lstdefinestyle{output}{
  basicstyle=\ttfamily\small,
  backgroundcolor=\color{gray!10},
  frame=single, rulecolor=\color{gray!50},
  breaklines=true,
}

\newtcolorbox{defbox}[2][]{colback=success_bg,colframe=green!50!black,
  fonttitle=\bfseries,title={Définition : #2},#1}
\newtcolorbox{theorembox}[2][]{colback=blue!5,colframe=suva_blue,
  fonttitle=\bfseries,title={Théorème : #2},#1}
\newtcolorbox{exercice}[2][]{colback=white,colframe=suva_blue!60,
  fonttitle=\bfseries,title={Exercice #2},#1,breakable}
\newtcolorbox{solution}[1][]{colback=gray!5,colframe=gray!40,
  fonttitle=\bfseries\itshape,title={Solution},#1,breakable}
\newtcolorbox{importantbox}[1][]{colback=suva_blue!10,colframe=suva_blue,
  fonttitle=\bfseries,title={Connexion avec le projet},#1}
\newtcolorbox{warningbox}[1][]{colback=warn_bg,colframe=orange!80!black,
  fonttitle=\bfseries,title={Attention},#1}

\newcommand{\Zq}{\mathbb{Z}_q}
\newcommand{\Rq}{R_q}
\newcommand{\ZZ}{\mathbb{Z}}
\newcommand{\pmod}[1]{\ (\mathrm{mod}\ #1)}

\newcommand{\auteur}{Mohamed Lamine MBENGUE}
\newcommand{\portfolio}{portfolio-alamine.web.app}

\pagestyle{fancy}
\fancyhf{}
\fancyhead[L]{\color{suva_blue}\textbf{SuVa --- Chapitre 2}}
\fancyhead[R]{\color{gray}Anneaux Polynomiaux}
\fancyfoot[L]{\footnotesize\color{gray}\auteur}
\fancyfoot[C]{\thepage}
\fancyfoot[R]{\footnotesize\color{gray}\portfolio}
\renewcommand{\headrulewidth}{0.4pt}
\renewcommand{\footrulewidth}{0.3pt}

% ══════════════════════════════════════════════════════════════════════
\begin{document}

\begin{center}
  {\color{suva_blue}\rule{\linewidth}{2pt}}\\[0.5cm]
  {\Huge\bfseries\color{suva_blue} Chapitre 2}\\[0.3cm]
  {\large\bfseries Anneaux Polynomiaux et $R_q = \mathbb{Z}_q[x]/(x^n+1)$}\\[0.2cm]
  {\small\color{gray} L'espace dans lequel vivent toutes les opérations de Dilithium}\\[0.3cm]
  {\color{suva_blue}\rule{\linewidth}{2pt}}
\end{center}

\begin{tcolorbox}[colback=suva_blue!5, colframe=suva_blue!40,
  left=6pt, right=6pt, top=4pt, bottom=4pt]
\small
\begin{tabular}{ll}
  \textbf{Auteur} & Mohamed Lamine MBENGUE \\
  \textbf{Spécialités} & Sécurité · DevOps · Dev Full Stack Web \& Mobile \\
  \textbf{Contact} & +221 78 178 44 36 · mbenguemohamedlamine2022@gmail.com \\
  \textbf{Portfolio} & \href{https://portfolio-alamine.web.app}{\texttt{portfolio-alamine.web.app}} \\
\end{tabular}
\end{tcolorbox}

\tableofcontents
\newpage

%══════════════════════════════════════════════════════════════════════
\section{Anneau de polynômes --- définition}
%══════════════════════════════════════════════════════════════════════

\subsection{Polynômes classiques}

Un polynôme à coefficients dans $\Zq$ est une expression~:
$$f(x) = f_0 + f_1 x + f_2 x^2 + \cdots + f_{n-1} x^{n-1}, \quad f_i \in \Zq$$

\textbf{Addition~:} $(f + g)_i = (f_i + g_i) \bmod q$

\textbf{Multiplication~:} $(f \cdot g)_k = \sum_{i+j=k} f_i g_j \bmod q$

La multiplication naïve (``schoolbook'') de deux polynômes de degré $n-1$
coûte $O(n^2)$ opérations. Pour $n = 256$, c'est $256^2 = 65\,536$
multiplications --- acceptable, mais lent.

\subsection{L'anneau quotient $R_q$}

\begin{defbox}{L'anneau $R_q$}
$$\Rq = \Zq[x] / (x^n + 1)$$
C'est l'ensemble des polynômes de degré $< n$ à coefficients dans $\Zq$,
où la multiplication est réduite modulo $x^n + 1$.

La réduction modulo $x^n + 1$ signifie que $x^n \equiv -1$.
\end{defbox}

\subsection{Pourquoi $x^n + 1$ et pas autre chose ?}

\begin{importantbox}
Dans le projet, $n = 256$ et le polynôme de réduction est $x^{256} + 1$.

Ce choix n'est pas anodin~:
\begin{enumerate}
  \item $x^{256} + 1$ est un \textbf{polynôme cyclotomique} ($x^{512} + 1$ factored).
  \item Il est \textbf{irréductible} sur $\Zq$ quand $q \equiv 1 \pmod{2n}$,
        ce que satisfait $q = 8\,380\,417$ (car $8\,380\,416 = 2^{13} \times 3 \times 256 + \ldots$).
  \item Ce choix permet l'utilisation de la \textbf{NTT} (Number Theoretic Transform)
        pour calculer la multiplication en $O(n \log n)$ au lieu de $O(n^2)$.
\end{enumerate}
\end{importantbox}

%══════════════════════════════════════════════════════════════════════
\section{Multiplication dans $R_q$ --- méthode naïve}
%══════════════════════════════════════════════════════════════════════

\subsection{Réduction modulo $x^n + 1$}

Quand on multiplie $f$ et $g$ et qu'on obtient un terme $x^k$ avec $k \geq n$~:
$$x^n \equiv -1 \Rightarrow x^{n+j} \equiv -x^j \Rightarrow x^{2n} \equiv 1$$

\textbf{Exemple avec $n = 4$~:}
$$\underbrace{x^5}_{\equiv -x} = x^{4+1} = x^4 \cdot x \equiv (-1) \cdot x = -x$$
$$\underbrace{x^7}_{\equiv -x^3} = x^{4+3} \equiv -x^3$$

\begin{lstlisting}[style=python]
def poly_mul_schoolbook(f, g, q, n=256):
    """
    Multiplication de polynomes dans R_q = Z_q[x]/(x^n + 1).
    Methode naïve O(n^2).
    """
    result = [0] * n
    for i, fi in enumerate(f):
        for j, gj in enumerate(g):
            k = i + j
            if k < n:
                result[k] = (result[k] + fi * gj) % q
            else:
                # x^(n+r) = -x^r  (car x^n = -1)
                result[k - n] = (result[k - n] - fi * gj) % q
    return result

# Test
q = 8_380_417
n = 4  # exemple avec n=4 pour lisibilite

f = [1, 2, 0, 0]  # f(x) = 1 + 2x
g = [0, 0, 1, 0]  # g(x) = x^2

# f * g = (1 + 2x) * x^2 = x^2 + 2x^3
result = poly_mul_schoolbook(f, g, q, n)
print(f"Resultat : {result}")  # [0, 0, 1, 2] = x^2 + 2x^3

# Test avec wrapping
f2 = [0, 0, 0, 1]  # f2(x) = x^3
g2 = [0, 0, 1, 0]  # g2(x) = x^2
# f2 * g2 = x^5 = -x  (car x^4 = -1 pour n=4)
result2 = poly_mul_schoolbook(f2, g2, q, n)
print(f"x^3 * x^2 = x^5 = -x : {result2}")
# Attendu : [0, q-1, 0, 0] = [0, 8380416, 0, 0]
\end{lstlisting}

%══════════════════════════════════════════════════════════════════════
\section{Implémentation dans le projet}
%══════════════════════════════════════════════════════════════════════

\subsection{La classe PolynomialRingDilithium}

Fichier~: \texttt{polynomials/polynomials\_generic.py}

\begin{lstlisting}[style=python]
# Voici ce que fait la classe (simplifie)
class PolynomialRingDilithium:
    def __init__(self, q=8_380_417, n=256):
        self.q = q
        self.n = n

    def __call__(self, coefficients):
        """Cree un polynome depuis une liste de coefficients."""
        return PolynomialDilithium(coefficients, self.q, self.n)

    def random_element(self):
        """Cree un polynome aleatoire."""
        import os
        coeffs = [int.from_bytes(os.urandom(3), 'little') % self.q
                  for _ in range(self.n)]
        return self(coeffs)

class PolynomialDilithium:
    def __init__(self, coeffs, q, n):
        self.coeffs = [c % q for c in coeffs]
        self.q = q
        self.n = n

    def __add__(self, other):
        """Addition dans R_q."""
        return PolynomialDilithium(
            [(a + b) % self.q for a, b in zip(self.coeffs, other.coeffs)],
            self.q, self.n
        )

    def __sub__(self, other):
        """Soustraction dans R_q."""
        return PolynomialDilithium(
            [(a - b) % self.q for a, b in zip(self.coeffs, other.coeffs)],
            self.q, self.n
        )

    def __mul__(self, other):
        """Multiplication dans R_q = Z_q[x]/(x^n+1)."""
        # Dans le projet reel, on utilise la NTT pour faire ca efficacement
        return poly_mul_schoolbook(self.coeffs, other.coeffs, self.q, self.n)

    def to_ntt(self):
        """Transforme en representation NTT (chapitre 3)."""
        ...  # utilise polyntt package

    def from_ntt(self):
        """Retourne depuis la representation NTT."""
        ...
\end{lstlisting}

\subsection{Opérations sur les modules (vecteurs et matrices)}

Fichier~: \texttt{modules/modules\_generic.py}

Dans Dilithium, on travaille avec des \textbf{vecteurs} et des \textbf{matrices}
de polynômes dans $\Rq$. Par exemple~:
\begin{itemize}
  \item $A \in \Rq^{k \times l}$ : matrice $4 \times 4$ de polynômes
  \item $s_1 \in \Rq^l$ : vecteur de $4$ polynômes
  \item $t = As_1 + s_2 \in \Rq^k$ : vecteur de $4$ polynômes
\end{itemize}

\begin{lstlisting}[style=python]
# Module = vecteur de polynomes
# Multiplication matrice-vecteur A @ s1

def matrix_vector_product(A, s1):
    """
    Calcule A * s1 dans R_q^k.
    A est k x l, s1 est de longueur l.
    Resultat : vecteur de longueur k.
    """
    k = len(A)      # nb de lignes
    l = len(s1)     # nb de colonnes
    result = []
    for i in range(k):
        row_product = sum(A[i][j] * s1[j] for j in range(l))
        # sum() d'elements de R_q : addition polynomiale
        result.append(row_product)
    return result
\end{lstlisting}

%══════════════════════════════════════════════════════════════════════
\section{Bit-packing --- sérialisation des polynômes}
%══════════════════════════════════════════════════════════════════════

Un coefficient de $t_1$ a besoin de $\lceil \log_2(q/2^d) \rceil$ bits.
Pour $q = 8\,380\,417$ et $d = 13$~: $q/2^{13} \approx 1023$, donc 10 bits.

La sérialisation compacte est essentielle pour réduire la taille de la
signature (2420 bytes pour Dilithium2).

\begin{lstlisting}[style=python]
def bit_pack_t1(t1_coeffs, n=256):
    """
    Pack les coefficients de t1 sur 10 bits chacun.
    t1 contient des valeurs dans [0, 1023].
    256 * 10 bits = 2560 bits = 320 bytes par polynome.
    """
    result = bytearray()
    for i in range(0, n, 4):
        # Emballer 4 coefficients de 10 bits = 5 bytes
        v0, v1, v2, v3 = t1_coeffs[i:i+4]
        # 10 bits par coeff, 4 coeffs = 40 bits = 5 bytes
        b0 = v0 & 0xFF
        b1 = ((v0 >> 8) & 0x03) | ((v1 & 0x3F) << 2)
        b2 = ((v1 >> 6) & 0x0F) | ((v2 & 0x0F) << 4)
        b3 = ((v2 >> 4) & 0x3F) | ((v3 & 0x03) << 6)
        b4 = (v3 >> 2) & 0xFF
        result.extend([b0, b1, b2, b3, b4])
    return bytes(result)

# Test
coeffs_t1 = list(range(256))  # valeurs 0, 1, ..., 255
packed = bit_pack_t1(coeffs_t1)
print(f"Taille packee : {len(packed)} bytes")  # 320 bytes
\end{lstlisting}

%══════════════════════════════════════════════════════════════════════
\section{Échantillonnage des polynômes secrets}
%══════════════════════════════════════════════════════════════════════

\subsection{sample\_in\_ball --- le polynôme challenge}

\begin{defbox}{sample\_in\_ball}
Crée un polynôme $c \in \Rq$ avec exactement $\tau = 39$ coefficients
égaux à $\pm 1$ et les autres à $0$.
Les positions et signes sont déterminés pseudoaléatoirement depuis
le hash $\tilde{c}$ (32 bytes).
\end{defbox}

\begin{lstlisting}[style=python]
def sample_in_ball(seed_bytes, tau, n=256, q=8_380_417):
    """
    Echantillonne un polynome challenge c avec tau coefficients +/-1.
    Algorithme : choisir tau positions aleatoires parmi n, puis signe.
    """
    import hashlib

    # Initialiser le XOF (SHAKE256)
    shake = hashlib.shake_256()
    shake.update(seed_bytes)

    coeffs = [0] * n

    # Remplir la "boule" de n a n-tau en arriere
    for i in range(n - tau, n):
        # Tirer j aleatoire dans [0, i]
        j = int.from_bytes(shake.digest(1), 'little') % (i + 1)
        # Placer le signe depuis le bit de poids fort
        sign_byte = int.from_bytes(shake.digest(1), 'little')
        sign = -1 if (sign_byte & 1) else 1
        # Echanger
        coeffs[i] = coeffs[j]
        coeffs[j] = sign

    return coeffs

# Verification : exactement tau = 39 coefficients non-nuls
c = sample_in_ball(b'\x00' * 32, tau=39)
nonzero = sum(1 for x in c if x != 0)
print(f"Coefficients non-nuls : {nonzero}")  # 39
print(f"Tous +1 ou -1 : {all(abs(x) <= 1 for x in c)}")
\end{lstlisting}

\subsection{rejection\_sample\_ntt\_poly --- la matrice A}

La matrice $A$ est générée par rejection sampling~: on génère des bytes
pseudo-aléatoires via SHAKE128 et on garde ceux qui sont dans $[0, q-1)$.

\begin{lstlisting}[style=python]
def rejection_sample(seed, q=8_380_417, n=256):
    """
    Genere un polynome a coefficients uniformes dans Z_q.
    Rejection sampling : rejette les valeurs >= q.
    """
    import hashlib
    coeffs = []
    shake = hashlib.shake_128()
    shake.update(seed)

    while len(coeffs) < n:
        # Lire 3 bytes -> 24 bits -> valeur dans [0, 2^24)
        raw = shake.digest(3)
        val = int.from_bytes(raw, 'little') & 0x7FFFFF  # 23 bits
        if val < q:
            coeffs.append(val)

    return coeffs[:n]

# Chaque appel avec le meme seed produit les memes coefficients
seed = b'\x01' + b'\x00' * 31  # 32 bytes
p1 = rejection_sample(seed)
p2 = rejection_sample(seed)
print(f"Deterministe ? {p1 == p2}")  # True
\end{lstlisting}

%══════════════════════════════════════════════════════════════════════
\section{Exercices}
%══════════════════════════════════════════════════════════════════════

\begin{exercice}{1 --- Multiplication dans $R_4$}
Soit $n = 4$, $q = 17$ et $R_q = \mathbb{Z}_{17}[x]/(x^4 + 1)$.
\begin{enumerate}[label=\alph*)]
  \item Calcule $(1 + x)(1 + x)$ dans $R_q$.
  \item Calcule $x^3 \cdot x^2$ dans $R_q$.
  \item Calcule $(2 + 3x^2)(1 + x)$ dans $R_q$.
\end{enumerate}
\end{exercice}

\begin{solution}
\begin{lstlisting}[style=python]
q, n = 17, 4

def poly_mul_small(f, g, q, n):
    res = [0] * n
    for i, fi in enumerate(f):
        for j, gj in enumerate(g):
            k = i + j
            if k < n: res[k] = (res[k] + fi * gj) % q
            else:     res[k-n] = (res[k-n] - fi * gj) % q
    return res

# a) (1+x)^2 dans Z_17[x]/(x^4+1)
f = [1, 1, 0, 0]
print(poly_mul_small(f, f, q, n))  # [1, 2, 1, 0] = 1 + 2x + x^2

# b) x^3 * x^2 = x^5 = x^4*x = -x
f2, g2 = [0,0,0,1], [0,0,1,0]
print(poly_mul_small(f2, g2, q, n))  # [0, 16, 0, 0] = -x (16 = -1 mod 17)

# c) (2 + 3x^2)(1 + x)
f3, g3 = [2,0,3,0], [1,1,0,0]
print(poly_mul_small(f3, g3, q, n))  # [2, 2, 3, 3]
\end{lstlisting}
\end{solution}

\begin{exercice}{2 --- Taille des clés}
Dans Dilithium2 ($n=256$, $q=8\,380\,417$, $k=l=4$, $d=13$)~:
\begin{enumerate}[label=\alph*)]
  \item La clé publique contient $\rho$ (32 bytes) + $t_1$ ($k$ polynômes, 10 bits/coeff).
        Calcule la taille totale de la clé publique.
  \item La signature contient $\tilde{c}$ (32 bytes) + $z$ ($l$ polynômes, $\lceil \log_2(2\gamma_1) \rceil$ bits/coeff)
        + $h$ (vecteur de bits compact). Estime la taille de la signature.
\end{enumerate}
\end{exercice}

\begin{solution}
\begin{lstlisting}[style=python]
import math
q = 8_380_417
n = 256
k = l = 4
d = 13
gamma_1 = 1 << 17  # 131072

# a) Taille cle publique
bits_t1_per_coeff = 10  # ceil(log2(q/2^d)) = ceil(log2(1023)) = 10
bytes_t1 = k * n * bits_t1_per_coeff // 8  # 4*256*10/8 = 1280 bytes
pk_size = 32 + bytes_t1
print(f"Cle publique : {pk_size} bytes")  # 1312 bytes

# b) Signature z
bits_z_per_coeff = math.ceil(math.log2(2 * gamma_1))  # = 18 bits
bytes_z = l * n * bits_z_per_coeff // 8  # 4*256*18/8 = 2304 bytes
# + c_tilde (32) + h (compact hint ~100 bytes)
# Total ~ 32 + 2304 + 84 = 2420 bytes
print(f"Signature z : ~{bytes_z} bytes (sans hint ni c_tilde)")
\end{lstlisting}
\end{solution}

\begin{exercice}{3 --- Exploration du code}
Ouvre \texttt{polynomials/polynomials.py} dans le projet.
\begin{enumerate}[label=\alph*)]
  \item Comment la multiplication \texttt{f * g} est-elle implémentée
        dans la vraie classe ? (Schoolbook ou NTT ?)
  \item Quelle méthode convertit un polynôme en représentation NTT ?
  \item Dans \texttt{Falcon\_dilithium.py}, trouve la ligne qui calcule
        $t = A \hat{s_1} + s_2$. Explique les opérations.
\end{enumerate}
\end{exercice}

\end{document}

\documentclass[12pt, a4paper]{article}
\usepackage[utf8]{inputenc}
\usepackage[T1]{fontenc}
\usepackage[french]{babel}
\usepackage[margin=2.5cm]{geometry}
\usepackage{amsmath, amssymb, amsthm, mathtools}
\usepackage{listings}
\usepackage{xcolor}
\usepackage{tcolorbox}
\tcbuselibrary{theorems, skins, breakable}
\usepackage{booktabs}
\usepackage{fancyhdr}
\usepackage{hyperref}
\usepackage{tikz}
\usetikzlibrary{arrows.meta, positioning, matrix}

\definecolor{suva_blue}{RGB}{30, 100, 180}
\definecolor{code_bg}{RGB}{245, 247, 250}
\definecolor{code_kw}{RGB}{0, 80, 200}
\definecolor{code_str}{RGB}{180, 40, 40}
\definecolor{code_cmt}{RGB}{80, 130, 80}
\definecolor{warn_bg}{RGB}{255, 248, 220}
\definecolor{success_bg}{RGB}{230, 255, 235}

\lstdefinestyle{python}{
  language=Python,
  basicstyle=\ttfamily\small,
  keywordstyle=\color{code_kw}\bfseries,
  stringstyle=\color{code_str},
  commentstyle=\color{code_cmt}\itshape,
  backgroundcolor=\color{code_bg},
  frame=single, rulecolor=\color{gray!40},
  numbers=left, numberstyle=\tiny\color{gray},
  breaklines=true, tabsize=4, showstringspaces=false,
}
\lstdefinestyle{output}{
  basicstyle=\ttfamily\small,
  backgroundcolor=\color{gray!10},
  frame=single, rulecolor=\color{gray!50},
  breaklines=true,
}

\newtcolorbox{defbox}[2][]{colback=success_bg,colframe=green!50!black,
  fonttitle=\bfseries,title={Définition : #2},#1}
\newtcolorbox{theorembox}[2][]{colback=blue!5,colframe=suva_blue,
  fonttitle=\bfseries,title={Théorème : #2},#1}
\newtcolorbox{exercice}[2][]{colback=white,colframe=suva_blue!60,
  fonttitle=\bfseries,title={Exercice #2},#1,breakable}
\newtcolorbox{solution}[1][]{colback=gray!5,colframe=gray!40,
  fonttitle=\bfseries\itshape,title={Solution},#1,breakable}
\newtcolorbox{importantbox}[1][]{colback=suva_blue!10,colframe=suva_blue,
  fonttitle=\bfseries,title={Connexion avec le projet},#1}
\newtcolorbox{warningbox}[1][]{colback=warn_bg,colframe=orange!80!black,
  fonttitle=\bfseries,title={Attention},#1}

\newcommand{\Zq}{\mathbb{Z}_q}
\newcommand{\Rq}{R_q}
\newcommand{\NTT}{\mathrm{NTT}}
\newcommand{\INTT}{\mathrm{NTT}^{-1}}

\pagestyle{fancy}
\fancyhf{}
\fancyhead[L]{\color{suva_blue}\textbf{SuVa --- Chapitre 3}}
\fancyhead[R]{\color{gray}NTT}
\fancyfoot[L]{\footnotesize\color{gray}Mohamed Lamine MBENGUE}
\fancyfoot[C]{\thepage}
\fancyfoot[R]{\footnotesize\color{gray}portfolio-alamine.web.app}
\renewcommand{\headrulewidth}{0.4pt}
\renewcommand{\footrulewidth}{0.3pt}

% ══════════════════════════════════════════════════════════════════════
\begin{document}

\begin{center}
  {\color{suva_blue}\rule{\linewidth}{2pt}}\\[0.5cm]
  {\Huge\bfseries\color{suva_blue} Chapitre 3}\\[0.3cm]
  {\large\bfseries Number Theoretic Transform (NTT)}\\[0.2cm]
  {\small\color{gray} Le moteur de performance de Dilithium}\\[0.3cm]
  {\color{suva_blue}\rule{\linewidth}{2pt}}
\end{center}

\begin{tcolorbox}[colback=suva_blue!5, colframe=suva_blue!40,
  left=6pt, right=6pt, top=4pt, bottom=4pt]
\small
\begin{tabular}{ll}
  \textbf{Auteur} & Mohamed Lamine MBENGUE \\
  \textbf{Spécialités} & Sécurité · DevOps · Dev Full Stack Web \& Mobile \\
  \textbf{Contact} & +221 78 178 44 36 · mbenguemohamedlamine2022@gmail.com \\
  \textbf{Portfolio} & \href{https://portfolio-alamine.web.app}{\texttt{portfolio-alamine.web.app}} \\
\end{tabular}
\end{tcolorbox}

\tableofcontents
\newpage

%══════════════════════════════════════════════════════════════════════
\section{Problème : multiplier des polynômes efficacement}
%══════════════════════════════════════════════════════════════════════

Dans Dilithium, on doit calculer $w = A \cdot y$ : une multiplication
matrice-vecteur de polynômes. Avec $k = l = 4$ et $n = 256$, c'est~:
$$4 \times 4 \times 256^2 = 1\,048\,576 \text{ multiplications en méthode naïve}$$

Avec la NTT, on réduit à~:
$$4 \times 4 \times (256 \log_2 256) = 4 \times 4 \times 2048 = 32\,768$$

Gain~: \textbf{facteur 32}.

%══════════════════════════════════════════════════════════════════════
\section{La DFT --- intuition}
%══════════════════════════════════════════════════════════════════════

La DFT (Discrete Fourier Transform) transforme un signal $f$ de longueur $n$~:
$$\hat{f}_k = \sum_{j=0}^{n-1} f_j \cdot \omega^{jk}, \quad k = 0, \ldots, n-1$$
où $\omega = e^{2\pi i / n}$ est la racine primitive $n$-ième de l'unité.

\textbf{Propriété clé~:} La multiplication devient une multiplication terme-à-terme~:
$$\widehat{f \cdot g}_k = \hat{f}_k \cdot \hat{g}_k$$

Donc~: $f \cdot g = \text{DFT}^{-1}(\text{DFT}(f) \odot \text{DFT}(g))$

\begin{itemize}
  \item DFT naïve~: $O(n^2)$
  \item FFT (Fast Fourier Transform)~: $O(n \log n)$
\end{itemize}

%══════════════════════════════════════════════════════════════════════
\section{La NTT --- version modulaire}
%══════════════════════════════════════════════════════════════════════

\begin{defbox}{NTT (Number Theoretic Transform)}
La NTT est la DFT sur $\Zq$ au lieu de $\mathbb{C}$.
Elle remplace $\omega = e^{2\pi i / n}$ par une \textbf{racine primitive
$n$-ième de l'unité dans $\Zq$}, notée $\zeta$.

$$\hat{f}_k = \sum_{j=0}^{n-1} f_j \cdot \zeta^{jk} \pmod{q}$$

\textbf{Condition~:} $\zeta$ existe dans $\Zq$ si $n \mid (q-1)$.
\end{defbox}

\subsection{Les racines de l'unité dans $\Zq$}

\begin{importantbox}
Pour Dilithium~: $q = 8\,380\,417$ et $n = 256$.

On a $q - 1 = 8\,380\,416 = 2^{23} \times \text{impair}$.
Donc $2n = 512$ divise $q - 1$ --- la NTT de taille $n = 256$ existe !

La racine primitive $2n = 512$-ième de l'unité est~:
$$\zeta_{512} = 1753 \pmod{q}$$
On vérifie~: $1753^{512} \equiv 1 \pmod{q}$ et $1753^{256} \not\equiv 1$.
\end{importantbox}

\begin{lstlisting}[style=python]
q = 8_380_417
zeta = 1753  # racine primitive 512-ieme

# Verification
print(f"zeta^512 = {pow(zeta, 512, q)}")  # 1
print(f"zeta^256 = {pow(zeta, 256, q)}")  # != 1 (devrait etre -1)
print(f"zeta^256 + 1 = {(pow(zeta, 256, q) + 1) % q}")  # 0

# Generer les puissances de zeta (twiddle factors)
zetas = [pow(zeta, k, q) for k in range(512)]
print(f"Les 8 premiers zeta^k : {zetas[:8]}")
\end{lstlisting}

%══════════════════════════════════════════════════════════════════════
\section{NTT pour $R_q = \Zq[x]/(x^n+1)$}
%══════════════════════════════════════════════════════════════════════

\subsection{NTT "twisted" (négacyclique)}

La NTT standard s'applique au produit cyclique $\Zq[x]/(x^n-1)$.
Pour $\Rq = \Zq[x]/(x^n+1)$, on utilise une version "twisted" (négacyclique).

\begin{theorembox}{NTT négacyclique}
Pour $f \in \Rq$ de taille $n$, la NTT négacyclique est~:
$$\hat{f}_k = \sum_{j=0}^{n-1} f_j \cdot \psi^j \cdot \omega^{jk}, \quad k = 0, \ldots, n-1$$
où~:
\begin{itemize}
  \item $\psi = \zeta_{2n}$ est la racine primitive $2n$-ième (ici~: $\zeta_{512} = 1753$)
  \item $\omega = \psi^2 = \zeta_n$ est la racine $n$-ième
\end{itemize}
\end{theorembox}

\textbf{Conséquence~:} La multiplication dans $\Rq$ devient~:
$$\text{NTT}(f \cdot g \bmod x^n+1) = \text{NTT}(f) \odot \text{NTT}(g)$$

Ce qui donne l'algorithme de multiplication efficace~:
\begin{enumerate}
  \item $\hat{f} = \NTT(f)$ \hfill$O(n \log n)$
  \item $\hat{g} = \NTT(g)$ \hfill$O(n \log n)$
  \item $\hat{h} = \hat{f} \odot \hat{g}$ (pointwise) \hfill$O(n)$
  \item $h = \INTT(\hat{h})$ \hfill$O(n \log n)$
\end{enumerate}

\subsection{Dans le projet}

\begin{importantbox}
Dans le code Python du projet, les polynômes ont deux représentations~:

\begin{itemize}
  \item \textbf{Représentation normale}~: \texttt{poly.coeffs} (liste de 256 entiers)
  \item \textbf{Représentation NTT}~: \texttt{poly\_hat.coeffs} (256 entiers en domaine NTT)
\end{itemize}

Les conversions se font avec~:
\begin{itemize}
  \item \texttt{poly.to\_ntt()} : transforme en NTT (forward transform)
  \item \texttt{poly\_hat.from\_ntt()} : NTT inverse
\end{itemize}

Quand on voit \texttt{A\_hat @ s1.to\_ntt()}, c'est~:
\begin{enumerate}
  \item Convertir $s_1$ en NTT
  \item Multiplier terme-à-terme chaque ligne de $\hat{A}$ avec $\hat{s_1}$
  \item Sommer les produits (dans le domaine NTT)
\end{enumerate}
\end{importantbox}

%══════════════════════════════════════════════════════════════════════
\section{Algorithme Cooley-Tukey --- implémentation}
%══════════════════════════════════════════════════════════════════════

\subsection{NTT itérative (butterfly)}

L'algorithme NTT opère par ``butterfly''~: à chaque étape, on divise
le problème en deux moitiés et on combine les résultats.

\begin{lstlisting}[style=python]
def ntt_forward(f, q=8_380_417, n=256, zeta=1753):
    """
    NTT negacyclique iterative (Cooley-Tukey).
    Transforme in-place le tableau f.
    """
    f = f.copy()
    psi = zeta  # racine 2n-ieme

    k = 1
    length = n // 2
    while length >= 1:
        for start in range(0, n, 2 * length):
            w = pow(psi, k, q)  # twiddle factor
            k += 1
            for j in range(start, start + length):
                t = (w * f[j + length]) % q
                f[j + length] = (f[j] - t) % q
                f[j] = (f[j] + t) % q
        length //= 2

    return f

def ntt_inverse(f_hat, q=8_380_417, n=256, zeta=1753):
    """
    NTT inverse.
    """
    f = f_hat.copy()
    psi = zeta
    psi_inv = pow(psi, q - 2, q)  # inverse de psi
    n_inv = pow(n, q - 2, q)       # inverse de n

    k = n - 1
    length = 1
    while length <= n // 2:
        for start in range(0, n, 2 * length):
            w = pow(psi_inv, k, q)
            k -= 1
            for j in range(start, start + length):
                t = f[j]
                f[j] = (t + f[j + length]) % q
                f[j + length] = (w * (f[j + length] - t)) % q
        length *= 2

    # Normalisation par n^(-1)
    for i in range(n):
        f[i] = (n_inv * f[i]) % q

    return f

# Test de coherence
import random
q = 8_380_417
n = 256

poly = [random.randint(0, q - 1) for _ in range(n)]
poly_hat = ntt_forward(poly, q, n)
poly_back = ntt_inverse(poly_hat, q, n)

print(f"NTT coherente : {poly == poly_back}")
\end{lstlisting}

\subsection{Multiplication via NTT}

\begin{lstlisting}[style=python]
def poly_mul_ntt(f, g, q=8_380_417, n=256):
    """
    Multiplication dans R_q via NTT.
    3x NTT(n log n) + 1x produit pointe (n).
    """
    f_hat = ntt_forward(f, q, n)
    g_hat = ntt_forward(g, q, n)

    # Produit pointe
    fg_hat = [(f_hat[i] * g_hat[i]) % q for i in range(n)]

    # NTT inverse
    return ntt_inverse(fg_hat, q, n)

# Verification : comparer avec schoolbook
import random
q, n = 8_380_417, 8  # petit n pour tester
f = [random.randint(0, q-1) for _ in range(n)]
g = [random.randint(0, q-1) for _ in range(n)]

result_ntt = poly_mul_ntt(f, g, q, n)
result_schoolbook = poly_mul_schoolbook(f, g, q, n)

print(f"NTT == Schoolbook : {result_ntt == result_schoolbook}")
\end{lstlisting}

%══════════════════════════════════════════════════════════════════════
\section{Les twiddle factors dans le projet Solidity}
%══════════════════════════════════════════════════════════════════════

\begin{importantbox}
Dans \texttt{ZKNOX\_NTT.sol}, les twiddle factors (puissances de $\zeta$)
sont \textbf{précompilés} et stockés dans des contrats déployés séparément.

À l'exécution, \texttt{extcodecopy} lit les constantes depuis ces contrats.
Cela évite de les recalculer on-chain à chaque vérification.

C'est le composant le plus gas-critique : écrit entièrement en \textbf{Yul}
(assembleur EVM) pour maximiser l'efficacité.
\end{importantbox}

Voici comment les twiddle factors sont générés en Python pour être
exportés vers Solidity~:

\begin{lstlisting}[style=python]
def generate_twiddle_factors(q=8_380_417, n=256, zeta=1753):
    """
    Genere les twiddle factors pour la NTT de taille n.
    Ces constantes sont hardcodees dans ZKNOX_NTT.sol.
    """
    psi = zeta  # racine 2n-ieme
    twiddles = []

    k = 1
    length = n // 2
    while length >= 1:
        for _ in range(n // (2 * length)):
            twiddles.append(pow(psi, k, q))
            k += 1
        length //= 2

    return twiddles

twiddles = generate_twiddle_factors()
print(f"Nombre de twiddle factors : {len(twiddles)}")  # 255

# Exporter en Solidity (format uint256[])
solidity_array = ", ".join(str(t) for t in twiddles)
print(f"uint256[255] constant TWIDDLES = [{solidity_array[:100]}...]")
\end{lstlisting}

%══════════════════════════════════════════════════════════════════════
\section{Exercices}
%══════════════════════════════════════════════════════════════════════

\begin{exercice}{1 --- Vérification de la racine de l'unité}
Soit $q = 8\,380\,417$ et $\zeta = 1753$.
\begin{enumerate}[label=\alph*)]
  \item Vérifie que $\zeta^{512} \equiv 1 \pmod{q}$.
  \item Vérifie que $\zeta^{256} \equiv -1 \pmod{q}$.
  \item Calcule $\zeta^{128}$ --- est-ce une racine $4$-ième de l'unité ?
  \item Pourquoi $\zeta^{256} = -1$ est-il lié à $x^{256} + 1 = 0$ ?
\end{enumerate}
\end{exercice}

\begin{solution}
\begin{lstlisting}[style=python]
q, zeta = 8_380_417, 1753

# a)
print(pow(zeta, 512, q))   # 1  -> zeta est bien racine 512-ieme

# b)
print(pow(zeta, 256, q))   # q-1 = 8380416 = -1 mod q

# c)
val = pow(zeta, 128, q)
print(f"zeta^128 = {val}")
print(f"(zeta^128)^4 = {pow(val, 4, q)}")   # 1 -> c'est bien une racine 4-ieme

# d) L'anneau Z_q[x]/(x^256+1) impose x^256 = -1.
# La NTT "voit" zeta comme x : zeta^256 = -1 correspond exactement
# a la relation x^256 = -1 dans l'anneau.
\end{lstlisting}
\end{solution}

\begin{exercice}{2 --- NTT de taille 4}
Soit $q = 17$, $n = 4$, et $\zeta_8 = 2$ (racine primitive $8$-ième modulo $17$).

Vérifie~: $2^8 \equiv 1 \pmod{17}$ et $2^4 \equiv -1 \pmod{17}$.

\begin{enumerate}[label=\alph*)]
  \item Calcule $\NTT([1, 2, 3, 4])$ manuellement.
  \item Calcule $\INTT$ du résultat et vérifie que tu retrouves $[1,2,3,4]$.
\end{enumerate}
\end{exercice}

\begin{solution}
\begin{lstlisting}[style=python]
q, n, zeta = 17, 4, 2

# Verification
print(pow(zeta, 8, q))   # 1
print(pow(zeta, 4, q))   # 16 = -1 mod 17

# NTT forward
f = [1, 2, 3, 4]
f_hat = ntt_forward(f, q, n, zeta)
print(f"NTT([1,2,3,4]) = {f_hat}")

# NTT inverse
f_back = ntt_inverse(f_hat, q, n, zeta)
print(f"INTT(NTT(f)) = {f_back}")   # doit etre [1,2,3,4]
\end{lstlisting}
\end{solution}

\begin{exercice}{3 --- Comparaison de performances}
\begin{enumerate}[label=\alph*)]
  \item En Python, mesure le temps de \texttt{poly\_mul\_schoolbook} vs
        \texttt{poly\_mul\_ntt} pour $n = 256$.
        (Utilise \texttt{import time; t = time.perf\_counter()})
  \item Quel gain de vitesse obtiens-tu ?
  \item Pourquoi le gain théorique ($n^2$ vs $n\log n$) et le gain
        mesuré peuvent différer ?
\end{enumerate}
\end{exercice}

\begin{solution}
\begin{lstlisting}[style=python]
import time, random

q, n = 8_380_417, 256
f = [random.randint(0, q-1) for _ in range(n)]
g = [random.randint(0, q-1) for _ in range(n)]

# Schoolbook
t0 = time.perf_counter()
for _ in range(10):
    poly_mul_schoolbook(f, g, q, n)
t_school = (time.perf_counter() - t0) / 10

# NTT
t0 = time.perf_counter()
for _ in range(10):
    poly_mul_ntt(f, g, q, n)
t_ntt = (time.perf_counter() - t0) / 10

print(f"Schoolbook : {t_school*1000:.2f} ms")
print(f"NTT        : {t_ntt*1000:.2f} ms")
print(f"Gain       : {t_school/t_ntt:.1f}x")

# Note : le gain mesure < n/log(n) car :
# - Overhead Python (interpretation)
# - pow() sur grands entiers est lent en Python
# Le vrai gain apparait en C/Rust/Yul
\end{lstlisting}
\end{solution}

\begin{exercice}{4 --- Le domaine NTT dans le code}
Dans \texttt{Falcon\_dilithium.py}, la ligne~:
\begin{center}
\texttt{t = (A\_hat @ s1.to\_ntt()).from\_ntt() + s2}
\end{center}
\begin{enumerate}[label=\alph*)]
  \item Pourquoi convertit-on $s_1$ avec \texttt{.to\_ntt()} avant de multiplier par $A$ ?
  \item $\hat{A}$ est-il déjà en représentation NTT ? Comment est-il généré ?
  \item Pourquoi appelle-t-on \texttt{.from\_ntt()} avant d'ajouter $s_2$ ?
  \item Quel serait le problème si on oubliait le \texttt{.from\_ntt()} ?
\end{enumerate}
\end{exercice}

\end{document}

\documentclass[12pt, a4paper]{article}
\usepackage[utf8]{inputenc}
\usepackage[T1]{fontenc}
\usepackage[french]{babel}
\usepackage[margin=2.5cm]{geometry}
\usepackage{amsmath, amssymb, amsthm, mathtools}
\usepackage{listings}
\usepackage{xcolor}
\usepackage{tcolorbox}
\tcbuselibrary{theorems, skins, breakable}
\usepackage{booktabs}
\usepackage{fancyhdr}
\usepackage{hyperref}
\usepackage{tikz}
\usetikzlibrary{arrows.meta, positioning}

\definecolor{suva_blue}{RGB}{30, 100, 180}
\definecolor{code_bg}{RGB}{245, 247, 250}
\definecolor{code_kw}{RGB}{0, 80, 200}
\definecolor{code_str}{RGB}{180, 40, 40}
\definecolor{code_cmt}{RGB}{80, 130, 80}
\definecolor{warn_bg}{RGB}{255, 248, 220}
\definecolor{success_bg}{RGB}{230, 255, 235}

\lstdefinestyle{python}{
  language=Python,
  basicstyle=\ttfamily\small,
  keywordstyle=\color{code_kw}\bfseries,
  stringstyle=\color{code_str},
  commentstyle=\color{code_cmt}\itshape,
  backgroundcolor=\color{code_bg},
  frame=single, rulecolor=\color{gray!40},
  numbers=left, numberstyle=\tiny\color{gray},
  breaklines=true, tabsize=4, showstringspaces=false,
}

\newtcolorbox{defbox}[2][]{colback=success_bg,colframe=green!50!black,
  fonttitle=\bfseries,title={Définition : #2},#1}
\newtcolorbox{theorembox}[2][]{colback=blue!5,colframe=suva_blue,
  fonttitle=\bfseries,title={Théorème : #2},#1}
\newtcolorbox{exercice}[2][]{colback=white,colframe=suva_blue!60,
  fonttitle=\bfseries,title={Exercice #2},#1,breakable}
\newtcolorbox{solution}[1][]{colback=gray!5,colframe=gray!40,
  fonttitle=\bfseries\itshape,title={Solution},#1,breakable}
\newtcolorbox{importantbox}[1][]{colback=suva_blue!10,colframe=suva_blue,
  fonttitle=\bfseries,title={Connexion avec le projet},#1}
\newtcolorbox{warningbox}[1][]{colback=warn_bg,colframe=orange!80!black,
  fonttitle=\bfseries,title={Attention},#1}

\pagestyle{fancy}
\fancyhf{}
\fancyhead[L]{\color{suva_blue}\textbf{SuVa --- Chapitre 4}}
\fancyhead[R]{\color{gray}Réseaux de Treillis}
\fancyfoot[L]{\footnotesize\color{gray}Mohamed Lamine MBENGUE}
\fancyfoot[C]{\thepage}
\fancyfoot[R]{\footnotesize\color{gray}portfolio-alamine.web.app}
\renewcommand{\headrulewidth}{0.4pt}
\renewcommand{\footrulewidth}{0.3pt}

% ══════════════════════════════════════════════════════════════════════
\begin{document}

\begin{center}
  {\color{suva_blue}\rule{\linewidth}{2pt}}\\[0.5cm]
  {\Huge\bfseries\color{suva_blue} Chapitre 4}\\[0.3cm]
  {\large\bfseries Réseaux de Treillis et Cryptographie Post-Quantique}\\[0.2cm]
  {\small\color{gray} La sécurité de Dilithium repose sur des problèmes difficiles sur les réseaux}\\[0.3cm]
  {\color{suva_blue}\rule{\linewidth}{2pt}}
\end{center}

\vspace{0.3cm}
\begin{tcolorbox}[colback=suva_blue!5, colframe=suva_blue!40,
  left=6pt, right=6pt, top=4pt, bottom=4pt]
\small
\begin{tabular}{ll}
  \textbf{Auteur}      & Mohamed Lamine MBENGUE \\
  \textbf{Spécialités} & Sécurité · DevOps · Dev Full Stack Web \& Mobile \\
  \textbf{Contact}     & +221 78 178 44 36 · mbenguemohamedlamine2022@gmail.com \\
  \textbf{Portfolio}   & \href{https://portfolio-alamine.web.app}{\texttt{portfolio-alamine.web.app}} \\
\end{tabular}
\end{tcolorbox}
\vspace{0.3cm}

\tableofcontents
\newpage

%══════════════════════════════════════════════════════════════════════
\section{Pourquoi la crypto actuelle sera cassée par les ordinateurs quantiques}
%══════════════════════════════════════════════════════════════════════

\subsection{ECDSA et la vulnérabilité quantique}

Ethereum utilise \textbf{ECDSA} (Elliptic Curve Digital Signature Algorithm).
Sa sécurité repose sur la difficulté du \textbf{problème du logarithme
discret sur les courbes elliptiques}~:

\begin{center}
Connaissant $P$ (point) et $Q = k \cdot P$, trouver $k$.
\end{center}

\textbf{L'algorithme de Shor} (1994) résout ce problème en temps polynomial
sur un ordinateur quantique. Conséquence~:
\begin{itemize}
  \item Un QC suffisant peut retrouver la clé privée depuis la clé publique
  \item Toutes les transactions Ethereum deviendraient forgeable
  \item $\sim$\$200 milliards de stablecoins exposés
\end{itemize}

\subsection{Pourquoi les réseaux résistent au quantique}

Les meilleurs algorithmes quantiques connus (Grover, Shor) n'offrent pas
d'avantage polynomial pour les problèmes sur les réseaux.

L'algorithme de Shor \textbf{ne fonctionne pas} sur les problèmes LWE/SIS.

%══════════════════════════════════════════════════════════════════════
\section{Qu'est-ce qu'un réseau de treillis (lattice) ?}
%══════════════════════════════════════════════════════════════════════

\begin{defbox}{Réseau de treillis}
Un réseau $\Lambda$ de rang $n$ est l'ensemble des combinaisons
\textbf{entières} de $n$ vecteurs linéairement indépendants $b_1, \ldots, b_n$~:
$$\Lambda = \left\{ \sum_{i=1}^n a_i b_i \mid a_i \in \mathbb{Z} \right\}$$
La matrice $B = [b_1 | \cdots | b_n]$ est appelée \textbf{base} du réseau.
\end{defbox}

\textbf{Visualisation 2D~:}

\begin{center}
\begin{tikzpicture}[scale=0.8]
  % Points du reseau
  \foreach \i in {-2,-1,0,1,2} {
    \foreach \j in {-2,-1,0,1,2} {
      \fill (\i*2 + \j*0.5, \j*1.5) circle (2pt);
    }
  }
  % Base vectors
  \draw[->, thick, suva_blue] (0,0) -- (2,0) node[right] {$b_1$};
  \draw[->, thick, red] (0,0) -- (0.5,1.5) node[above] {$b_2$};
  % Axes
  \draw[gray, ->] (-4,0) -- (5.5,0);
  \draw[gray, ->] (0,-3) -- (0,4);
\end{tikzpicture}
\end{center}

%══════════════════════════════════════════════════════════════════════
\section{Problèmes difficiles sur les réseaux}
%══════════════════════════════════════════════════════════════════════

\subsection{SVP --- Shortest Vector Problem}

\begin{defbox}{SVP}
\textbf{Problème du vecteur le plus court~:}
Trouver le vecteur non-nul le plus court dans un réseau $\Lambda$.
$$\lambda_1(\Lambda) = \min_{v \in \Lambda \setminus \{0\}} \|v\|$$
\end{defbox}

SVP est NP-difficile dans le cas général. Les meilleurs algorithmes
(BKZ, LLL) ont une complexité exponentielle.

\subsection{LWE --- Learning With Errors}

\begin{defbox}{LWE}
\textbf{Problème LWE~:} Connaissant les couples $(a_i, b_i)$ où~:
$$b_i = \langle a_i, s \rangle + e_i \pmod{q}$$
avec $s$ le \textbf{secret}, $a_i$ aléatoires uniformes, et $e_i$ un
\textbf{bruit} petit, retrouver $s$.
\end{defbox}

\begin{importantbox}
\textbf{Analogie~:} Imagine qu'on te donne 100 équations du type~:
$$a_1 x_1 + a_2 x_2 + \cdots + a_n x_n + e \equiv b \pmod{q}$$
Chaque équation a une petite erreur $e$. Sans les erreurs, le système
est facilement résoluble (Gauss). \textbf{Avec les erreurs}, même avec
des millions d'équations, retrouver $(x_1, \ldots, x_n)$ est
computationnellement infaisable.
\end{importantbox}

\subsection{MLWE --- Module Learning With Errors}

\begin{defbox}{MLWE}
Version modulaire de LWE~: le secret $s$ et le bruit $e$ sont des
\textbf{vecteurs de polynômes} dans $R_q^l$, et $A$ est une matrice
$k \times l$ de polynômes dans $R_q$.

$$t = A s_1 + s_2 \pmod{q}$$

où $s_1 \in R_q^l$ et $s_2 \in R_q^k$ ont des coefficients petits
(bornés par $\eta$).
\end{defbox}

\begin{importantbox}
Dans Dilithium~:
\begin{itemize}
  \item $A \in R_q^{4 \times 4}$ est publique (générée de $\rho$)
  \item $s_1 \in R_q^4$ a des coefficients dans $[-\eta, \eta] = [-2, 2]$ (secret)
  \item $s_2 \in R_q^4$ a des coefficients dans $[-\eta, \eta] = [-2, 2]$ (secret)
  \item $t = As_1 + s_2$ est publique (dans la clé publique)
\end{itemize}
\textbf{Sécurité~:} Étant donné $(A, t)$, retrouver $(s_1, s_2)$ est MLWE.
\end{importantbox}

\subsection{MSIS --- Module Short Integer Solution}

\begin{defbox}{MSIS}
\textbf{Module Short Integer Solution~:}
Étant donné $A \in R_q^{k \times l}$, trouver un vecteur court non-nul
$z \in R_q^l$ tel que~:
$$Az \equiv 0 \pmod{q}$$
\end{defbox}

\begin{importantbox}
\textbf{Lien avec Dilithium~:} L'unforgeability (impossibilité de forger
une signature) repose sur MSIS. Un attaquant qui forge une signature
produit essentiellement un court vecteur $z$ tel que $Az \approx ct_1 \cdot 2^d$.
Cela serait une solution au problème MSIS.
\end{importantbox}

%══════════════════════════════════════════════════════════════════════
\section{L'attaque BKZ --- comment mesurer la sécurité}
%══════════════════════════════════════════════════════════════════════

\begin{defbox}{Algorithme BKZ}
BKZ (Block Korkine-Zolotarev) est le meilleur algorithme pratique
pour réduire une base de réseau (trouver de courts vecteurs).
Il paramétrisé par un \textbf{block size} $\beta$~:

\begin{itemize}
  \item Plus $\beta$ est grand, plus le résultat est bon
  \item Mais le temps d'exécution est $\sim 2^{0.292\beta}$ (classique)
  \item Et $\sim 2^{0.265\beta}$ (quantique, algorithme de Grover)
\end{itemize}

La sécurité en bits est $\approx 0.292\beta$.
\end{defbox}

\subsection{Estimation de sécurité pour Dilithium2}

Les scripts dans \texttt{security-estimates-simon/} calculent le
$\beta$ nécessaire pour casser le schéma~:

\begin{lstlisting}[style=python]
# Simplifie depuis Falcon_MLWE_security.py
def estimate_mlwe_security(n, q, k, l, eta):
    """
    Estime la securite de l'instance MLWE.
    Trouve le block-size BKZ minimal pour casser le schema.
    """
    # Dimension du reseau : (k+l)*n
    lattice_dim = (k + l) * n

    # Cherche le beta optimal
    for beta in range(50, 1000):
        # Longueur du vecteur cible (approximation)
        # Pour trouver s1, s2 avec coefficients <= eta
        target_norm = eta * (lattice_dim ** 0.5)

        # Longueur minimale que BKZ peut trouver (Hermite factor)
        hermite_factor = (
            (beta / (2 * 3.14159 * 2.71828)) **
            (1 / (2 * beta)) *
            (q ** (k * n / lattice_dim))
        )
        expected_norm = hermite_factor ** lattice_dim

        if expected_norm <= target_norm:
            # Ce beta suffit pour l'attaque
            security_classical = 0.292 * beta
            security_quantum = 0.265 * beta
            return beta, security_classical, security_quantum

    return None

# Dilithium2
q = 8_380_417
n = 256; k = 4; l = 4; eta = 2
beta, sec_c, sec_q = estimate_mlwe_security(n, q, k, l, eta)
print(f"Block size BKZ : {beta}")
print(f"Securite classique : {sec_c:.0f} bits")
print(f"Securite quantique : {sec_q:.0f} bits")
# Attendu : ~426 bits, ~124 bits classique, ~112 bits quantique
\end{lstlisting}

%══════════════════════════════════════════════════════════════════════
\section{SelfTargetMSIS --- sécurité de la signature}
%══════════════════════════════════════════════════════════════════════

La sécurité de Dilithium repose sur un problème encore plus fort~:
le \textbf{SelfTargetMSIS}. Forger une signature sans la clé secrète
revient à trouver un vecteur court $(z, c)$ tel que~:
$$Az - ct_1 \cdot 2^d \approx 0 \pmod{q}$$
avec $\|z\|_\infty < \gamma_1 - \beta$ et $c$ un polynôme sparse.

C'est encore plus difficile que MSIS car le target $ct_1 \cdot 2^d$ dépend
du message à signer.

%══════════════════════════════════════════════════════════════════════
\section{Exercices}
%══════════════════════════════════════════════════════════════════════

\begin{exercice}{1 --- LWE jouet}
Soit $q = 101$, $n = 3$, secret $s = (2, 7, 1)$, bruit max $|e| \leq 1$.

On te donne ces échantillons LWE~:
\begin{align*}
  (a_1, b_1) &= ((3, 1, 4), 3 \cdot 2 + 1 \cdot 7 + 4 \cdot 1 + e_1)\\
  (a_2, b_2) &= ((1, 5, 9), 1 \cdot 2 + 5 \cdot 7 + 9 \cdot 1 + e_2)
\end{align*}

\begin{enumerate}[label=\alph*)]
  \item Calcule les $b_i$ exacts (sans erreur) en arithmétique $\mathbb{Z}_{101}$.
  \item Ajoute un bruit $e_1 = 1$, $e_2 = -1$ et calcule les $b_i$ bruitées.
  \item Si tu as accès à plein d'échantillons propres (sans bruit), peux-tu
        récupérer $s$ facilement ? (Indice~: Gauss)
  \item Avec le bruit, pourquoi est-ce difficile ?
\end{enumerate}
\end{exercice}

\begin{solution}
\begin{lstlisting}[style=python]
q = 101
s = [2, 7, 1]

a1, a2 = [3, 1, 4], [1, 5, 9]

# a) Sans bruit
b1_exact = sum(a*si for a, si in zip(a1, s)) % q
b2_exact = sum(a*si for a, si in zip(a2, s)) % q
print(f"b1 exact = {b1_exact}")   # (6+7+4) = 17
print(f"b2 exact = {b2_exact}")   # (2+35+9) = 46

# b) Avec bruit
b1_noisy = (b1_exact + 1) % q
b2_noisy = (b2_exact - 1) % q
print(f"b1 bruite = {b1_noisy}")   # 18
print(f"b2 bruite = {b2_noisy}")   # 45

# c) Sans bruit : on peut faire Gauss.
# Avec seulement 3 variables et 3 equations propres :
# 3x + y + 4z = 17
# x + 5y + 9z = 46
# ... -> solution unique = (2, 7, 1)

# d) Avec bruit : b = As + e.
# On ne sait pas quels echantillons ont quel bruit.
# Gauss donne une solution fausse.
# Les meilleurs algorithmes prennent un temps exponentiel.
\end{lstlisting}
\end{solution}

\begin{exercice}{2 --- MLWE dans Dilithium}
Dans Dilithium2 ($k=l=4$, $n=256$, $q=8\,380\,417$, $\eta=2$)~:
\begin{enumerate}[label=\alph*)]
  \item Quelle est la dimension totale du réseau utilisé pour l'attaque ?
  \item Si on augmentait $\eta$ de 2 à 4, que se passe-t-il pour la sécurité ?
  \item Si on augmentait $q$ (corps plus grand), que se passe-t-il ?
\end{enumerate}
\end{exercice}

\begin{solution}
\begin{lstlisting}[style=python]
# a) Dimension du reseau
n, k, l = 256, 4, 4
dim = (k + l) * n
print(f"Dimension : {dim}")  # 2048

# b) eta plus grand => vecteur secret plus long => plus facile a trouver
# => securite DIMINUE si on augmente eta
# C'est pourquoi Dilithium3 augmente k, l plutot qu'eta.

# c) q plus grand => le reseau est "plus facile" a reduire proportionnellement
# => securite DIMINUE si on augmente q avec les memes k, l.
# C'est le probleme du ZK-Dilithium avec BabyBear (q=2^31-1) :
# plus grand q => il faut augmenter k, l pour compenser.
\end{lstlisting}
\end{solution}

\begin{exercice}{3 --- Explorer les scripts de sécurité}
Ouvre \texttt{security-estimates-simon/Falcon\_Dilithium.py}.
\begin{enumerate}[label=\alph*)]
  \item Quelle est la sécurité estimée pour Dilithium2 en bits (classique et quantique) ?
  \item Comment la sécurité est-elle estimée (quelle méthode, quel modèle) ?
  \item Quelle différence entre attaque MLWE et MSIS dans le contexte de Dilithium ?
\end{enumerate}
\end{exercice}

\end{document}

\documentclass[12pt, a4paper]{article}
\usepackage[utf8]{inputenc}
\usepackage[T1]{fontenc}
\usepackage[french]{babel}
\usepackage[margin=2.5cm]{geometry}
\usepackage{amsmath, amssymb, amsthm, mathtools}
\usepackage{listings}
\usepackage{xcolor}
\usepackage{tcolorbox}
\tcbuselibrary{theorems, skins, breakable}
\usepackage{booktabs}
\usepackage{fancyhdr}
\usepackage{hyperref}
\usepackage{tikz}
\usetikzlibrary{arrows.meta, positioning, shapes.geometric}
\usepackage{algorithm}
\usepackage{algpseudocode}

\definecolor{suva_blue}{RGB}{30, 100, 180}
\definecolor{code_bg}{RGB}{245, 247, 250}
\definecolor{code_kw}{RGB}{0, 80, 200}
\definecolor{code_str}{RGB}{180, 40, 40}
\definecolor{code_cmt}{RGB}{80, 130, 80}
\definecolor{warn_bg}{RGB}{255, 248, 220}
\definecolor{success_bg}{RGB}{230, 255, 235}

\lstdefinestyle{python}{
  language=Python,
  basicstyle=\ttfamily\small,
  keywordstyle=\color{code_kw}\bfseries,
  stringstyle=\color{code_str},
  commentstyle=\color{code_cmt}\itshape,
  backgroundcolor=\color{code_bg},
  frame=single, rulecolor=\color{gray!40},
  numbers=left, numberstyle=\tiny\color{gray},
  breaklines=true, tabsize=4, showstringspaces=false,
}
\lstdefinestyle{output}{
  basicstyle=\ttfamily\small,
  backgroundcolor=\color{gray!10},
  frame=single, rulecolor=\color{gray!50},
  breaklines=true,
}

\newtcolorbox{defbox}[2][]{colback=success_bg,colframe=green!50!black,
  fonttitle=\bfseries,title={Définition : #2},#1}
\newtcolorbox{theorembox}[2][]{colback=blue!5,colframe=suva_blue,
  fonttitle=\bfseries,title={Théorème : #2},#1}
\newtcolorbox{exercice}[2][]{colback=white,colframe=suva_blue!60,
  fonttitle=\bfseries,title={Exercice #2},#1,breakable}
\newtcolorbox{solution}[1][]{colback=gray!5,colframe=gray!40,
  fonttitle=\bfseries\itshape,title={Solution},#1,breakable}
\newtcolorbox{importantbox}[1][]{colback=suva_blue!10,colframe=suva_blue,
  fonttitle=\bfseries,title={Connexion avec le projet},#1}
\newtcolorbox{warningbox}[1][]{colback=warn_bg,colframe=orange!80!black,
  fonttitle=\bfseries,title={Attention},#1}

\newcommand{\Zq}{\mathbb{Z}_q}
\newcommand{\Rq}{R_q}

\pagestyle{fancy}
\fancyhf{}
\fancyhead[L]{\color{suva_blue}\textbf{SuVa --- Chapitre 5}}
\fancyhead[R]{\color{gray}Dilithium Complet}
\fancyfoot[L]{\footnotesize\color{gray}Mohamed Lamine MBENGUE}
\fancyfoot[C]{\thepage}
\fancyfoot[R]{\footnotesize\color{gray}portfolio-alamine.web.app}
\renewcommand{\headrulewidth}{0.4pt}
\renewcommand{\footrulewidth}{0.3pt}

% ══════════════════════════════════════════════════════════════════════
\begin{document}

\begin{center}
  {\color{suva_blue}\rule{\linewidth}{2pt}}\\[0.5cm]
  {\Huge\bfseries\color{suva_blue} Chapitre 5}\\[0.3cm]
  {\large\bfseries Dilithium / ML-DSA --- Le schéma de signature complet}\\[0.2cm]
  {\small\color{gray} Keygen, Sign, Verify : code source et mathématiques détaillés}\\[0.3cm]
  {\color{suva_blue}\rule{\linewidth}{2pt}}
\end{center}

\vspace{0.3cm}
\begin{tcolorbox}[colback=suva_blue!5, colframe=suva_blue!40,
  left=6pt, right=6pt, top=4pt, bottom=4pt]
\small
\begin{tabular}{ll}
  \textbf{Auteur}      & Mohamed Lamine MBENGUE \\
  \textbf{Spécialités} & Sécurité · DevOps · Dev Full Stack Web \& Mobile \\
  \textbf{Contact}     & +221 78 178 44 36 · mbenguemohamedlamine2022@gmail.com \\
  \textbf{Portfolio}   & \href{https://portfolio-alamine.web.app}{\texttt{portfolio-alamine.web.app}} \\
\end{tabular}
\end{tcolorbox}
\vspace{0.3cm}

\tableofcontents
\newpage

%══════════════════════════════════════════════════════════════════════
\section{Vue d'ensemble du schéma Dilithium}
%══════════════════════════════════════════════════════════════════════

\subsection{Paramètres Dilithium2 (niveau utilisé dans SuVa)}

\begin{center}
\begin{tabular}{clp{7cm}}
\toprule
\textbf{Param.} & \textbf{Valeur} & \textbf{Rôle} \\
\midrule
$n$    & 256 & Degré des polynômes \\
$q$    & 8\,380\,417 & Module premier \\
$k$    & 4 & Lignes de la matrice $A$ \\
$l$    & 4 & Colonnes de la matrice $A$ \\
$\eta$ & 2 & Borne des coefficients secrets ($s_1, s_2$) \\
$\tau$ & 39 & Nb de $\pm 1$ dans le challenge $c$ \\
$\gamma_1$ & $2^{17} = 131\,072$ & Plage de masquage de $z$ \\
$\gamma_2$ & $(q-1)/88 = 95\,232$ & Plage de rounding \\
$d$    & 13 & Bits droppés de la clé publique \\
$\omega$ & 80 & Max de 1s dans le hint $h$ \\
$\beta$ & 78 & Borne de rejet ($= \tau \times \eta$) \\
\bottomrule
\end{tabular}
\end{center}

\subsection{Les trois algorithmes}

\begin{center}
\begin{tikzpicture}[
  block/.style={draw=suva_blue, rounded corners, fill=blue!5,
    minimum width=5cm, minimum height=1.2cm,
    text width=4.8cm, align=center},
  arrow/.style={->, thick, suva_blue},
  label/.style={font=\small, color=gray}
]
  \node[block] (kg) {
    \textbf{KeyGen} (off-chain, une fois)\\
    \footnotesize $\zeta \to (A, s_1, s_2) \to t = As_1+s_2$\\
    \footnotesize \texttt{pk = (rho, t1), sk = (rho, K, tr, s1, s2, t0)}
  };
  \node[block, below=1.5cm of kg] (sg) {
    \textbf{Sign} (off-chain, par message)\\
    \footnotesize Boucle de rejet jusqu'à trouver $(z, h)$ valide\\
    \footnotesize \texttt{sig = (c\_tilde, z, h)}
  };
  \node[block, below=1.5cm of sg] (vf) {
    \textbf{Verify} (on-chain Solidity)\\
    \footnotesize Recalcule $w' = Az - ct_1 \cdot 2^d$\\
    \footnotesize Compare $\tilde{c} \stackrel{?}{=} H(\mu \| w')$
  };
  \draw[arrow] (kg) -- node[label, right] {sk → Sign} (sg);
  \draw[arrow] (sg) -- node[label, right] {(pk, sig, msg)} (vf);
  \node[right=2cm of kg, font=\scriptsize, align=left, color=gray] {
    Fichier :\\
    \texttt{Falcon\_dilithium.py}
  };
  \node[right=2cm of vf, font=\scriptsize, align=left, color=gray] {
    Fichier :\\
    \texttt{ZKNOX\_dilithium.sol}
  };
\end{tikzpicture}
\end{center}

%══════════════════════════════════════════════════════════════════════
\section{KeyGen --- Génération de clés en détail}
%══════════════════════════════════════════════════════════════════════

\subsection{Algorithme}

\begin{lstlisting}[style=python]
# Extrait de Falcon_dilithium/Falcon_dilithium.py - methode keygen()

def keygen(self, _xof=None, _xof2=None):
    # 1. Graine principale : 32 bytes aleatoires
    zeta = self.random_bytes(32)

    # 2. Expander : SHAKE256(zeta) -> 128 bytes
    #    Divise en rho (32), rho_prime (64), K (32)
    seed_bytes = self._h(zeta, 128)
    rho        = seed_bytes[:32]   # seed pour matrice publique A
    rho_prime  = seed_bytes[32:96] # seed pour vecteurs secrets s1, s2
    K          = seed_bytes[96:]   # seed pour la signature (nonce)

    # 3. Generer la matrice publique A in R_q^{k x l}
    # A est generee par rejection sampling depuis SHAKE128(rho)
    # IMPORTANTE : A est dans le domaine NTT deja (A_hat)
    A_hat = self._expand_matrix_from_seed(rho)
    # A_hat est k x l = 4x4 matrices de polynomes

    # 4. Generer les vecteurs secrets
    # s1 in R_q^l : coefficients dans [-eta, eta]
    # s2 in R_q^k : coefficients dans [-eta, eta]
    s1, s2 = self._expand_vector_from_seed(rho_prime)
    # Pour Dilithium2 (eta=2) : coefficients dans {-2,-1,0,1,2}

    # 5. Calculer t = A * s1 + s2
    t = (A_hat @ s1.to_ntt()).from_ntt() + s2
    # Decomposition : t = t1 * 2^d + t0
    t1, t0 = t.power_2_round(self.d)

    # 6. Construire les cles
    pk = self._pack_pk(rho, t1)    # 1312 bytes
    tr = self._h(pk, 32)           # Hash de la cle publique
    sk = self._pack_sk(rho, K, tr, s1, s2, t0)  # 2528 bytes

    return pk, sk
\end{lstlisting}

\subsection{Tailles}

\begin{center}
\begin{tabular}{lll}
\toprule
\textbf{Composant} & \textbf{Taille} & \textbf{Contenu} \\
\midrule
$\rho$ & 32 bytes & Graine pour générer $A$ \\
$t_1$ & 1280 bytes & $k \times n \times 10$ bits $= 4 \times 256 \times 10 / 8$ \\
\textbf{Total pk} & \textbf{1312 bytes} & \\
\midrule
$\rho$ & 32 bytes & \\
$K$ & 32 bytes & \\
$\text{tr}$ & 32 bytes & Hash de pk \\
$s_1$ & 384 bytes & $l \times n \times (2\eta+1)$-bits codage \\
$s_2$ & 384 bytes & $k \times n \times (2\eta+1)$-bits codage \\
$t_0$ & 1664 bytes & $k \times n \times 13$ bits \\
\textbf{Total sk} & \textbf{2528 bytes} & \\
\bottomrule
\end{tabular}
\end{center}

%══════════════════════════════════════════════════════════════════════
\section{Sign --- La boucle de rejet}
%══════════════════════════════════════════════════════════════════════

\begin{warningbox}
La boucle de rejet est le cœur de Dilithium.
En moyenne, \textbf{4.25 itérations} sont nécessaires pour Dilithium2.
\end{warningbox}

\begin{lstlisting}[style=python]
# Extrait simplifie de sign()

def sign(self, sk_bytes, m, _xof=None, _xof2=None):
    rho, K, tr, s1, s2, t0 = self._unpack_sk(sk_bytes)
    A_hat = self._expand_matrix_from_seed(rho)

    # Message representative : mu = SHAKE256(tr || m, 64 bytes)
    mu = self._h(tr + m, 64)

    # rho_prime pour le masque y
    rho_prime = self._h(K + mu, 64)

    # Precomputer en NTT (important pour la performance)
    s1 = s1.to_ntt()
    s2 = s2.to_ntt()
    t0 = t0.to_ntt()

    alpha = 2 * self.gamma_2
    kappa = 0  # nonce incremente a chaque iteration

    while True:  # Boucle de rejet
        # Etape 1 : Tirer le masque y
        # y in R_q^l avec ||y||_inf <= gamma_1
        y = self._expand_mask_vector(rho_prime, kappa)
        kappa += self.l

        # Etape 2 : Calculer w = A*y dans R_q^k
        w = (A_hat @ y.to_ntt()).from_ntt()

        # Etape 3 : Decomposer w en w1 (haute) et w0 (basse)
        w1, w0 = w.decompose(alpha)

        # Etape 4 : Hacher pour obtenir le challenge
        # c_tilde = H(mu || w1_packed)
        c_tilde = self._h(mu + w1.bit_pack_w(self.gamma_2), 32)
        c = self.R.sample_in_ball(c_tilde, self.tau).to_ntt()

        # Etape 5 : z = y + c*s1
        z = y + (s1.scale(c)).from_ntt()

        # REJET 1 : z trop grand -> fuit info sur s1
        if z.check_norm_bound(self.gamma_1 - self.beta):
            continue  # Recommencer

        # Etape 6 : Calculer les termes de correction
        w0_minus_cs2 = w0 - s2.scale(c).from_ntt()

        # REJET 2 : w0 - cs2 trop grand
        if w0_minus_cs2.check_norm_bound(self.gamma_2 - self.beta):
            continue

        c_t0 = t0.scale(c).from_ntt()

        # REJET 3 : c*t0 trop grand
        if c_t0.check_norm_bound(self.gamma_2):
            continue

        # Etape 7 : Construire le hint
        h = (w0_minus_cs2 + c_t0).make_hint_optimised(w1, alpha)

        # REJET 4 : trop de 1s dans h (trop grande correction)
        if h.sum_hint() > self.omega:
            continue

        # Succes ! Empaqueter la signature
        return self._pack_sig(c_tilde, z, h)
        # Signature = (c_tilde, z, h) = (32, ~1152, ~100) bytes
\end{lstlisting}

\subsection{Pourquoi le rejet protège la clé secrète ?}

\begin{theorembox}{Rejet pour la confidentialité}
Si on publiait $z = y + cs_1$ sans restriction, un attaquant qui observe
plusieurs signatures pourrait faire de l'analyse statistique~:

$z = y + cs_1 \Rightarrow z - y = cs_1$

Si la distribution de $z$ dépend de $s_1$, la clé secrète fuit.

\textbf{Le rejet uniformise~:} En rejetant les $z$ trop grands (proches
de la borne), on s'assure que la distribution de $z$ est essentiellement
\textbf{uniforme} sur $[-\gamma_1+\beta, \gamma_1-\beta]$ et n'informe
pas sur $s_1$.
\end{theorembox}

%══════════════════════════════════════════════════════════════════════
\section{Verify --- Vérification complète}
%══════════════════════════════════════════════════════════════════════

\begin{lstlisting}[style=python]
# Extrait simplifie de verify()

def verify(self, pk_bytes, m, sig_bytes, _xof=None):
    # Decompresser
    rho, t1 = self._unpack_pk(pk_bytes)
    c_tilde, z, h = self._unpack_sig(sig_bytes)

    # Verifications basiques
    if h.sum_hint() > self.omega:
        return False
    if z.check_norm_bound(self.gamma_1 - self.beta):
        return False

    # Regenerer A (depuis la graine rho)
    A_hat = self._expand_matrix_from_seed(rho)

    # Recalculer mu
    tr = self._h(pk_bytes, 32)
    mu = self._h(tr + m, 64)

    # Reconstruire c depuis c_tilde
    c = self.R.sample_in_ball(c_tilde, self.tau).to_ntt()

    # Calculer Az - c*t1*2^d
    # t1 * 2^d reconstruit le vecteur t1 * 2^d dans Z_q
    z_ntt = z.to_ntt()
    t1_scaled = t1.scale(1 << self.d).to_ntt()
    Az_minus_ct1 = ((A_hat @ z_ntt) - t1_scaled.scale(c)).from_ntt()
    # Note : (A*z - c*t1*2^d) ≈ A*y (si z = y + c*s1 et t1*2^d ≈ t-t0)

    # Utiliser le hint pour recuperer w1'
    w_prime = h.use_hint(Az_minus_ct1, 2 * self.gamma_2)

    # Verification finale
    c_tilde_check = self._h(mu + w_prime.bit_pack_w(self.gamma_2), 32)
    return c_tilde == c_tilde_check
\end{lstlisting}

\subsection{Preuve de correction (intuition)}

Lors de la signature on a~:
$$z = y + cs_1, \quad t = As_1 + s_2, \quad t = t_1 \cdot 2^d + t_0$$

Lors de la vérification~:
\begin{align*}
Az - ct_1 \cdot 2^d
&= A(y + cs_1) - ct_1 \cdot 2^d \\
&= Ay + cAs_1 - ct_1 \cdot 2^d \\
&= Ay + c(t - s_2) - c(t - t_0) \\
&= Ay - cs_2 + ct_0 \\
&\approx Ay \quad \text{(si } \|cs_2\|_\infty, \|ct_0\|_\infty < \gamma_2\text{)}
\end{align*}

Le hint corrige la petite différence entre $Ay$ et $Az - ct_1 \cdot 2^d$
pour retrouver $w_1 = \text{HighBits}(Ay)$.

%══════════════════════════════════════════════════════════════════════
\section{ETHDilithium --- les deux changements clés}
%══════════════════════════════════════════════════════════════════════

\subsection{Changement 1 : SHAKE → Keccak256}

\begin{center}
\begin{tabular}{lll}
\toprule
& \textbf{Dilithium standard} & \textbf{ETHDilithium} \\
\midrule
Hash function & SHAKE-128/256 & Keccak-256 \\
Implémentation Solidity & 269 lignes & opcode natif (3 gaz) \\
Coût on-chain & $\sim 50\%$ du gas total & $\sim 5\%$ \\
Conformité NIST & Oui (FIPS 204) & Non \\
\bottomrule
\end{tabular}
\end{center}

\subsection{Changement 2 : Clé publique étendue}

\begin{center}
\begin{tabular}{lll}
\toprule
& \textbf{Dilithium} & \textbf{ETHDilithium} \\
\midrule
Clé publique & $(\rho, t_1)$ & $(\hat{A}, t_1)$ \\
Taille pk & 1312 bytes & $\sim 4096 + 1280$ bytes \\
Coût keygen & Normal & On précalcule $\hat{A}$ \\
Coût verify & Doit régénérer $\hat{A}$ & $\hat{A}$ déjà là \\
Économie gas & --- & $\sim -50\%$ \\
\bottomrule
\end{tabular}
\end{center}

\begin{lstlisting}[style=python]
# Extrait de Falcon_dilithium_eth.py

def keygen(self, _xof=None, _xof2=None):
    # ... (pareil jusqu'a t1, t0)

    # DIFFERENCE : stocker A_hat dans la cle publique
    # Au lieu de (rho, t1), on stocke (A_hat, t1_new)
    pk_new = self._pack_pk_new(A_hat, t1)
    # pk_new est plus grande (toute la matrice A)
    # mais la verification n'a plus besoin de regenerer A

    # t1_new : deja en NTT form dans la pk
    # Pour que la verification puisse directement utiliser
    # t1 * 2^d sans recalculer le NTT
    return pk_new, sk
\end{lstlisting}

%══════════════════════════════════════════════════════════════════════
\section{Exercices}
%══════════════════════════════════════════════════════════════════════

\begin{exercice}{1 --- Tracer le flow complet}
Trace sur papier le flow complet de Dilithium en annotant chaque étape
avec le fichier/fonction correspondant dans le projet~:

\begin{enumerate}[label=\alph*)]
  \item KeyGen~: de \texttt{random\_bytes(32)} jusqu'à \texttt{pack\_pk}
  \item Sign~: de \texttt{\_unpack\_sk} jusqu'à \texttt{return \_pack\_sig}
  \item Verify~: de \texttt{\_unpack\_pk} jusqu'au \texttt{return True/False}
\end{enumerate}
\end{exercice}

\begin{exercice}{2 --- Pourquoi le nonce kappa}
Dans \texttt{sign()}, \texttt{kappa} est incrémenté de $l$ à chaque itération.
\begin{enumerate}[label=\alph*)]
  \item Pourquoi ne pas utiliser \texttt{os.urandom()} pour le masque $y$ ?
  \item Que se passerait-il si deux signatures utilisaient le même $y$ ?
  \item Comment \texttt{kappa} garantit-il que $y$ est différent à chaque itération ?
\end{enumerate}
\end{exercice}

\begin{solution}
\begin{lstlisting}[style=python]
# b) Si y est reutilise pour deux messages m1, m2 :
# sig1 = (c1, z1=y+c1*s1, h1)
# sig2 = (c2, z2=y+c2*s1, h2)
# Alors : z1 - z2 = (c1 - c2) * s1
# Si c1 != c2, on peut isoler s1 ! -> CATASTROPHIQUE

# c) kappa est incremente : rho_prime || kappa
# -> deux iterations differentes -> y differents
# Le PRNG est deterministe depuis (rho_prime, kappa)
# -> securite ET determinisme (utile pour les KAT)
\end{lstlisting}
\end{solution}

\begin{exercice}{3 --- Sécurité du hint}
Le hint $h$ peut contenir au plus $\omega = 80$ valeurs 1.
\begin{enumerate}[label=\alph*)]
  \item Combien de bits d'information le hint encode-t-il au maximum ?
        ($k \times n$ bits potentiels, mais $\leq \omega$ sont 1)
  \item Pourquoi la limite $\omega$ est-elle nécessaire pour la sécurité ?
  \item Que se passe-t-il si un attaquant peut forcer $\|h\|_1 > \omega$ ?
\end{enumerate}
\end{exercice}

\begin{exercice}{4 --- ETHDilithium vs Dilithium en Python}
\begin{enumerate}[label=\alph*)]
  \item Ouvre \texttt{example.py} et \texttt{example\_zk.py}.
  \item Lance les deux exemples et observe la sortie.
  \item Mesure et compare les temps d'exécution de ETHDilithium2 vs Dilithium2
        pour keygen, sign et verify.
  \item Pourquoi ETHDilithium est plus lent en Python mais plus rapide on-chain ?
\end{enumerate}
\end{exercice}

\end{document}

\documentclass[12pt, a4paper]{article}
\usepackage[utf8]{inputenc}
\usepackage[T1]{fontenc}
\usepackage[french]{babel}
\usepackage[margin=2.5cm]{geometry}
\usepackage{amsmath, amssymb}
\usepackage{listings}
\usepackage{xcolor}
\usepackage{tcolorbox}
\tcbuselibrary{theorems, skins, breakable}
\usepackage{booktabs}
\usepackage{fancyhdr}
\usepackage{hyperref}
\usepackage{tikz}
\usetikzlibrary{arrows.meta, positioning, shapes.geometric}

\definecolor{suva_blue}{RGB}{30, 100, 180}
\definecolor{code_bg}{RGB}{245, 247, 250}
\definecolor{code_kw}{RGB}{0, 80, 200}
\definecolor{code_str}{RGB}{180, 40, 40}
\definecolor{code_cmt}{RGB}{80, 130, 80}
\definecolor{warn_bg}{RGB}{255, 248, 220}
\definecolor{success_bg}{RGB}{230, 255, 235}
\definecolor{sol_green}{RGB}{0, 120, 80}
\definecolor{sol_bg}{RGB}{240, 255, 245}

\lstdefinestyle{solidity}{
  language=C++,
  basicstyle=\ttfamily\small,
  keywordstyle=\color{code_kw}\bfseries,
  stringstyle=\color{code_str},
  commentstyle=\color{code_cmt}\itshape,
  backgroundcolor=\color{sol_bg},
  frame=single, rulecolor=\color{sol_green!40},
  numbers=left, numberstyle=\tiny\color{gray},
  breaklines=true, tabsize=4, showstringspaces=false,
  morekeywords={uint256, uint8, bytes, bytes32, bool, address, memory,
    storage, calldata, external, view, pure, returns, mapping, struct,
    event, emit, require, revert, modifier, payable, public, private,
    internal, virtual, override, assembly, let, mstore, mload},
}
\lstdefinestyle{python}{
  language=Python,
  basicstyle=\ttfamily\small,
  keywordstyle=\color{code_kw}\bfseries,
  stringstyle=\color{code_str},
  commentstyle=\color{code_cmt}\itshape,
  backgroundcolor=\color{code_bg},
  frame=single, rulecolor=\color{gray!40},
  numbers=left, numberstyle=\tiny\color{gray},
  breaklines=true, tabsize=4, showstringspaces=false,
}
\lstdefinestyle{bash}{
  basicstyle=\ttfamily\small\color{white},
  backgroundcolor=\color{gray!80},
  frame=single, rulecolor=\color{gray!80},
  breaklines=true,
}
\lstdefinestyle{output}{
  basicstyle=\ttfamily\small,
  backgroundcolor=\color{gray!10},
  frame=single, rulecolor=\color{gray!50},
  breaklines=true,
}

\newtcolorbox{defbox}[2][]{colback=success_bg,colframe=green!50!black,
  fonttitle=\bfseries,title={Définition : #2},#1}
\newtcolorbox{exercice}[2][]{colback=white,colframe=suva_blue!60,
  fonttitle=\bfseries,title={Exercice #2},#1,breakable}
\newtcolorbox{solution}[1][]{colback=gray!5,colframe=gray!40,
  fonttitle=\bfseries\itshape,title={Solution},#1,breakable}
\newtcolorbox{importantbox}[1][]{colback=suva_blue!10,colframe=suva_blue,
  fonttitle=\bfseries,title={Connexion avec le projet},#1}
\newtcolorbox{warningbox}[1][]{colback=warn_bg,colframe=orange!80!black,
  fonttitle=\bfseries,title={Attention},#1}
\newtcolorbox{gasbox}[1][]{colback=orange!5,colframe=orange!80,
  fonttitle=\bfseries,title={Gas cost},#1}

\pagestyle{fancy}
\fancyhf{}
\fancyhead[L]{\color{suva_blue}\textbf{SuVa --- Chapitre 6}}
\fancyhead[R]{\color{gray}Ethereum / Solidity / ERC-4337}
\fancyfoot[L]{\footnotesize\color{gray}Mohamed Lamine MBENGUE}
\fancyfoot[C]{\thepage}
\fancyfoot[R]{\footnotesize\color{gray}portfolio-alamine.web.app}
\renewcommand{\headrulewidth}{0.4pt}
\renewcommand{\footrulewidth}{0.3pt}

% ══════════════════════════════════════════════════════════════════════
\begin{document}

\begin{center}
  {\color{suva_blue}\rule{\linewidth}{2pt}}\\[0.5cm]
  {\Huge\bfseries\color{suva_blue} Chapitre 6}\\[0.3cm]
  {\large\bfseries Ethereum, Solidity, Gas et ERC-4337}\\[0.2cm]
  {\small\color{gray} La partie on-chain du projet SuVa}\\[0.3cm]
  {\color{suva_blue}\rule{\linewidth}{2pt}}
\end{center}

\vspace{0.3cm}
\begin{tcolorbox}[colback=suva_blue!5, colframe=suva_blue!40,
  left=6pt, right=6pt, top=4pt, bottom=4pt]
\small
\begin{tabular}{ll}
  \textbf{Auteur}      & Mohamed Lamine MBENGUE \\
  \textbf{Spécialités} & Sécurité · DevOps · Dev Full Stack Web \& Mobile \\
  \textbf{Contact}     & +221 78 178 44 36 · mbenguemohamedlamine2022@gmail.com \\
  \textbf{Portfolio}   & \href{https://portfolio-alamine.web.app}{\texttt{portfolio-alamine.web.app}} \\
\end{tabular}
\end{tcolorbox}
\vspace{0.3cm}

\tableofcontents
\newpage

%══════════════════════════════════════════════════════════════════════
\section{L'EVM --- Ethereum Virtual Machine}
%══════════════════════════════════════════════════════════════════════

\subsection{Fonctionnement de base}

L'EVM est une machine à pile (\emph{stack machine}) qui exécute le bytecode
des smart contracts. Chaque opération (opcode) a un coût en \textbf{gas}.

\begin{defbox}{Gas}
Le \textbf{gas} est l'unité de mesure du coût de calcul sur Ethereum.
\begin{itemize}
  \item Chaque opcode coûte un certain nombre de gas
  \item Le gas est payé en ETH (gas\_used $\times$ gas\_price)
  \item La limite de gas par bloc est $\sim 30\,000\,000$ gas
\end{itemize}
\end{defbox}

\subsection{Opcodes importants pour ce projet}

\begin{center}
\begin{tabular}{llr}
\toprule
\textbf{Opcode} & \textbf{Opération} & \textbf{Gas} \\
\midrule
\texttt{ADD} & Addition 256-bit & 3 \\
\texttt{MUL} & Multiplication 256-bit & 5 \\
\texttt{DIV} & Division 256-bit & 5 \\
\texttt{ADDMOD} & (a+b) mod n & 8 \\
\texttt{MULMOD} & (a*b) mod n & 8 \\
\texttt{KECCAK256} & Hash Keccak-256 & 30 + 6/word \\
\texttt{MLOAD} & Lecture mémoire & 3 \\
\texttt{MSTORE} & Écriture mémoire & 3 \\
\texttt{EXTCODECOPY} & Copier bytecode & 100 + copie \\
\bottomrule
\end{tabular}
\end{center}

\begin{gasbox}
Comparaison des coûts de vérification post-quantique~:
\begin{center}
\begin{tabular}{lr}
\toprule
\textbf{Schéma} & \textbf{Gas} \\
\midrule
ECDSA (actuel Ethereum) & $\sim$ 3\,000 \\
\textbf{Falcon (NIST FIPS 206)} & $\sim$ \textbf{350\,000} \\
ETHDilithium (Keccak) & $\sim$ 6\,600\,000 \\
Dilithium standard (SHAKE) & $\sim$ 13\,500\,000 \\
SPHINCS+ & $\sim$ 50--500\,000\,000 \\
Limite par bloc & 30\,000\,000 \\
\bottomrule
\end{tabular}
\end{center}
Falcon est $\sim 20\times$ moins cher que Dilithium et tient bien dans la limite de bloc.
\end{gasbox}

%══════════════════════════════════════════════════════════════════════
\section{Solidity --- Bases essentielles}
%══════════════════════════════════════════════════════════════════════

\subsection{Types de données utiles dans ce projet}

\begin{lstlisting}[style=solidity]
// Types de base
uint256 a = 42;              // entier 256-bit non signe
bytes32 h;                   // 32 bytes fixes (hash)
bytes memory data;           // tableau de bytes dynamique
bool ok = true;

// Tableaux
uint256[] memory arr;                 // tableau dynamique
uint256[256] memory poly;             // polynome (256 coefficients)
uint256[][] memory vector;            // vecteur de polynomes
uint256[][][] memory matrix;          // matrice de polynomes

// Structures (utilises pour PubKey et Signature)
struct PubKey {
    bytes32 rho;            // graine de la matrice A
    bytes tr;               // hash de la cle publique
    uint256[][] t1;         // vecteur t1 (k polynomes)
}

struct Signature {
    bytes32[] c_tilde;      // hash du challenge (32 bytes)
    uint256[][] z;          // vecteur z (l polynomes)
    uint256[][] h;          // hint (k polynomes de bits)
}
\end{lstlisting}

\subsection{Les zones mémoire}

\begin{defbox}{Memory vs Storage vs Calldata}
\begin{itemize}
  \item \textbf{storage}~: donnees persistantes sur la blockchain (très cher)
  \item \textbf{memory}~: donnees temporaires pendant l'execution (moderé)
  \item \textbf{calldata}~: arguments en lecture seule d'un appel (le moins cher)
\end{itemize}
\end{defbox}

\begin{importantbox}
Dans les contrats Dilithium, presque tout est en \texttt{memory}
car la vérification est stateless (pas de stockage permanent).
Chaque appel à \texttt{verify()} recommence de zéro.
\end{importantbox}

\subsection{Fonctions et visibilité}

\begin{lstlisting}[style=solidity]
// Fonction de verification (interface externe)
function verify(
    PubKey memory pk,
    bytes memory message,
    Signature memory signature
) external view returns (bool) {
    // external : appelee depuis l'exterieur
    // view     : ne modifie pas l'etat
    // returns  : type de retour
    ...
}

// Fonction interne (helper)
function compute_c_ntt(
    Signature memory sig
) internal pure returns (uint256[] memory) {
    // internal : uniquement depuis ce contrat
    // pure     : pas de lecture d'etat non plus
    ...
}
\end{lstlisting}

%══════════════════════════════════════════════════════════════════════
\section{Yul --- EVM Assembly dans Solidity}
%══════════════════════════════════════════════════════════════════════

\subsection{Pourquoi Yul dans ce projet ?}

Le NTT (\texttt{ZKNOX\_NTT.sol}) est écrit entièrement en Yul pour
maximiser l'efficacité du gas. En Solidity pur, le compilateur ajoute
des vérifications et des copies inutiles. Yul donne un contrôle total.

\subsection{Syntaxe de base Yul}

\begin{lstlisting}[style=solidity]
// Dans Solidity, on utilise assembly { ... } pour du Yul inline
function mulmod_example(uint256 a, uint256 b, uint256 m)
    internal pure returns (uint256 result) {
    assembly {
        // Variables locales
        let x := mulmod(a, b, m)   // (a * b) mod m

        // Operations sur la memoire
        let ptr := mload(0x40)     // Lire le free memory pointer
        mstore(ptr, x)             // Ecrire x a l'adresse ptr

        // Affectation du resultat
        result := x
    }
}

// NTT butterfly step (simplifie)
function butterfly(uint256 u, uint256 v, uint256 w, uint256 q)
    internal pure returns (uint256 u_out, uint256 v_out) {
    assembly {
        // t = w * v mod q
        let t := mulmod(w, v, q)
        // u_out = (u + t) mod q
        u_out := addmod(u, t, q)
        // v_out = (u - t) mod q
        v_out := addmod(u, sub(q, t), q)
    }
}
\end{lstlisting}

\subsection{extcodecopy --- lire les twiddle factors}

\begin{lstlisting}[style=solidity]
// Dans ZKNOX_NTT.sol
// Les twiddle factors sont stockes comme bytecode d'un contrat
// (deploye par ZKNOX_dilithium_deploy.sol)

function load_twiddle_factors(address twiddle_contract)
    internal view returns (uint256[] memory) {
    assembly {
        // Taille du bytecode (= nb_twiddles * 32 bytes)
        let size := extcodesize(twiddle_contract)

        // Allouer la memoire
        let ptr := mload(0x40)
        mstore(0x40, add(ptr, size))

        // Copier le bytecode (qui contient les constantes)
        extcodecopy(twiddle_contract, ptr, 0, size)
    }
}
// Avantage : extcodecopy est moins cher que lire depuis storage
// Inconvenient : les constantes doivent etre deployees d'abord
\end{lstlisting}

%══════════════════════════════════════════════════════════════════════
\section{Architecture des contrats Solidity dans le projet}
%══════════════════════════════════════════════════════════════════════

\begin{center}
\begin{tikzpicture}[
  block/.style={draw=sol_green, rounded corners, fill=sol_bg,
    minimum width=4cm, minimum height=0.8cm, text width=3.8cm,
    align=center, font=\small},
  arrow/.style={->, thick, sol_green!80}
]
  \node[block] (entry) {\texttt{ZKNOX\_dilithium.sol}\\{\tiny entry point, verify()}};
  \node[block, below left=1cm and 0.5cm of entry]
    (core) {\texttt{ZKNOX\_dilithium\_core.sol}\\{\tiny core\_1(), core\_2()}};
  \node[block, below right=1cm and 0.5cm of entry]
    (eth) {\texttt{ZKNOX\_ethdilithium.sol}\\{\tiny variant Keccak}};
  \node[block, below=1cm of core]
    (ntt) {\texttt{ZKNOX\_NTT.sol}\\{\tiny NTT en Yul}};
  \node[block, below left=1cm and 0.5cm of ntt]
    (nttdil) {\texttt{ZKNOX\_NTT\_dilithium.sol}\\{\tiny NTTFW, NTTINV}};
  \node[block, below right=1cm and 0.5cm of ntt]
    (utils) {\texttt{ZKNOX\_dilithium\_utils.sol}\\{\tiny constantes, structs}};
  \node[block, below=1cm of core, xshift=3.5cm]
    (hint) {\texttt{ZKNOX\_hint.sol}\\{\tiny useHint()}};
  \node[block, right=1cm of entry]
    (shake) {\texttt{ZKNOX\_shake.sol}\\{\tiny SHAKE-256 Solidity}};
  \node[block, below=0.5cm of shake]
    (keccak) {\texttt{ZKNOX\_keccak\_prng.sol}\\{\tiny Keccak PRNG}};

  \draw[arrow] (entry) -- (core);
  \draw[arrow] (entry) -- (eth);
  \draw[arrow] (core) -- (ntt);
  \draw[arrow] (ntt) -- (nttdil);
  \draw[arrow] (ntt) -- (utils);
  \draw[arrow] (core) -- (hint);
  \draw[arrow] (entry) -- (shake);
  \draw[arrow] (eth) -- (keccak);
\end{tikzpicture}
\end{center}

\subsection{ZKNOX\_dilithium.sol --- le contrat principal}

\begin{lstlisting}[style=solidity]
// ZKNOX_dilithium.sol (simplifie)
pragma solidity ^0.8.25;

import "./ZKNOX_dilithium_core.sol";
import "./ZKNOX_shake.sol";
import "./ZKNOX_NTT_dilithium.sol";

contract ZKNOX_dilithium is ZKNOX_dilithium_core {
    // Adresses des contrats de twiddle factors (deployes)
    address immutable apsirev;    // NTT forward
    address immutable apsiInvrev; // NTT inverse

    constructor(address _apsirev, address _apsiInvrev) {
        apsirev = _apsirev;
        apsiInvrev = _apsiInvrev;
    }

    function verify(
        PubKey memory pk,
        bytes memory msgs,
        Signature memory signature
    ) external view returns (bool) {
        // STEP 1: Valider la structure de la signature
        (uint256 norm_h, uint256[][] memory h, uint256[][] memory z) =
            dilithium_core_1(signature);

        if (norm_h > OMEGA) return false;

        // Verifier que z est dans [-gamma_1+beta, gamma_1-beta]
        for (uint i = 0; i < L; i++) {
            for (uint j = 0; j < N; j++) {
                if (z[i][j] > GAMMA_1_MINUS_BETA &&
                    Q - z[i][j] > GAMMA_1_MINUS_BETA) {
                    return false;
                }
            }
        }

        // STEP 2: Reconstruire c en domaine NTT
        uint256[] memory c_ntt = compute_c_ntt(signature);

        // STEP 3: t1 * 2^d en NTT
        uint256[][] memory t1_new = ZKNOX_Expand_Vec(pk.t1);
        for (uint i = 0; i < K; i++) {
            for (uint j = 0; j < N; j++) {
                t1_new[i][j] <<= D;  // shift par 2^d
            }
            t1_new[i] = ZKNOX_NTTFW(t1_new[i], apsirev);
        }

        // STEP 4: w' = UseHint(h, A*z - c*t1)
        bytes memory w_prime = dilithium_core_2(
            apsirev, apsiInvrev, pk, z, c_ntt, h, t1_new
        );

        // STEP 5: Verification du hash final
        ctx_shake memory ctx;
        ctx = shake_update(ctx, pk.tr);
        ctx = shake_update(ctx, msgs);
        bytes memory mu = shake_digest(ctx, 64);

        ctx_shake memory ctx2;
        ctx2 = shake_update(ctx2, mu);
        ctx2 = shake_update(ctx2, w_prime);
        bytes32 final_hash = bytes32(shake_digest(ctx2, 32));

        for (uint i = 0; i < 32; i++) {
            if (signature.c_tilde[i] != uint8(final_hash[i]))
                return false;
        }
        return true;
    }
}
\end{lstlisting}

%══════════════════════════════════════════════════════════════════════
\section{ERC-4337 --- Account Abstraction}
%══════════════════════════════════════════════════════════════════════

\subsection{Pourquoi ERC-4337 ?}

Ethereum valide nativement les transactions via ECDSA.
Pour utiliser une signature post-quantique, on ne peut pas changer
le protocole Ethereum. ERC-4337 permet de contourner ça via des
\textbf{smart contract wallets}.

\subsection{Composants ERC-4337}

\begin{center}
\begin{tikzpicture}[
  block/.style={draw=suva_blue, rounded corners, fill=blue!5,
    minimum width=3.5cm, minimum height=1cm, text width=3.3cm,
    align=center, font=\small},
  arrow/.style={->, thick, suva_blue}
]
  \node[block] (user) {Utilisateur\\{\tiny Génère UserOperation}};
  \node[block, right=1.5cm of user] (mempool)
    {UserOperation\\Mempool\\{\tiny File d'attente}};
  \node[block, right=1.5cm of mempool]
    (bundler) {Bundler\\{\tiny Agrège les ops}};
  \node[block, below=1.5cm of bundler]
    (entry) {EntryPoint\\Contract\\{\tiny Contrat central}};
  \node[block, below=1.5cm of entry]
    (wallet) {Smart Wallet\\(SuVa)\\{\tiny validateUserOp()}};
  \node[block, right=1.5cm of entry]
    (paymaster) {Paymaster\\{\tiny Paie le gas}};

  \draw[arrow] (user) -- (mempool);
  \draw[arrow] (mempool) -- (bundler);
  \draw[arrow] (bundler) -- (entry);
  \draw[arrow] (entry) -- (wallet);
  \draw[arrow] (entry) -- (paymaster);
\end{tikzpicture}
\end{center}

\subsection{Le Smart Wallet SuVa}

\begin{lstlisting}[style=solidity]
// Contrat SuVa Wallet (a construire !)
// Ce contrat sera implemente dans le futur du projet

pragma solidity ^0.8.25;

import "@account-abstraction/contracts/interfaces/IAccount.sol";
import "./ZKNOX_ethdilithium.sol";

contract SuVaWallet is IAccount {
    // Cle publique post-quantique de l'utilisateur
    PubKey public quantum_pk;

    // Verifieur ETHDilithium
    ZKNOX_ethdilithium immutable verifier;

    constructor(
        PubKey memory _quantum_pk,
        address _verifier
    ) {
        quantum_pk = _quantum_pk;
        verifier = ZKNOX_ethdilithium(_verifier);
    }

    // Fonction appelee par EntryPoint pour valider une transaction
    function validateUserOp(
        UserOperation calldata userOp,
        bytes32 userOpHash,
        uint256 missingAccountFunds
    ) external returns (uint256 validationData) {
        // La signature dans userOp contient la signature post-quantique
        Signature memory sig = abi.decode(
            userOp.signature,
            (Signature)
        );

        // Verifier avec ETHDilithium
        bool valid = verifier.verify(
            quantum_pk,
            abi.encodePacked(userOpHash),
            sig
        );

        // 0 = valide, 1 = invalide
        return valid ? 0 : 1;
    }
}
\end{lstlisting}

\begin{importantbox}
Le flow complet pour une transaction SuVa~:
\begin{enumerate}
  \item L'utilisateur génère une signature Dilithium/Falcon off-chain (Python/SDK)
  \item La signature est encodée dans \texttt{UserOperation.signature}
  \item Le Bundler soumet la \texttt{UserOperation} à l'EntryPoint
  \item L'EntryPoint appelle \texttt{validateUserOp()} du Smart Wallet
  \item Le Smart Wallet appelle \texttt{verify()} du contrat ZKNOX
  \item Si la vérification passe, la transaction est exécutée
\end{enumerate}
\end{importantbox}

%══════════════════════════════════════════════════════════════════════
\section{Foundry --- Tester les contrats}
%══════════════════════════════════════════════════════════════════════

\subsection{Structure d'un test Foundry}

\begin{lstlisting}[style=solidity]
// test/ZKNOX_dilithium.t.sol (genere par Python)
pragma solidity ^0.8.25;

import "forge-std/Test.sol";
import "../src/ZKNOX_dilithium.sol";

contract TestDilithium is Test {
    ZKNOX_dilithium verifier;

    // Donnees de test (generees par generate_dilithium_test_vectors.py)
    bytes constant PUBLIC_KEY = hex"abc123..."; // 1312 bytes
    bytes constant MESSAGE = "We are ZKNox.";
    bytes constant SIGNATURE = hex"def456..."; // 2420 bytes

    function setUp() public {
        // Deployer le contrat (et les twiddle factors)
        // ...
        verifier = new ZKNOX_dilithium(...);
    }

    function test_verify_valid_signature() public {
        // Decoder les donnees
        PubKey memory pk = abi.decode(PUBLIC_KEY, (PubKey));
        Signature memory sig = abi.decode(SIGNATURE, (Signature));

        // Verifier
        bool result = verifier.verify(pk, MESSAGE, sig);
        assertTrue(result, "La signature valide doit passer");
    }

    function test_verify_wrong_message() public {
        PubKey memory pk = abi.decode(PUBLIC_KEY, (PubKey));
        Signature memory sig = abi.decode(SIGNATURE, (Signature));

        bool result = verifier.verify(pk, "Wrong message", sig);
        assertFalse(result, "Mauvais message doit echouer");
    }

    function test_gas_cost() public {
        // Foundry mesure automatiquement le gas si on utilise
        // vm.startSnapshotGas / vm.stopSnapshotGas
        uint256 gasBefore = gasleft();
        // ... (verification)
        uint256 gasUsed = gasBefore - gasleft();
        emit log_uint(gasUsed);  // ~6.6M pour ETHDilithium
    }
}
\end{lstlisting}

\subsection{Commandes Foundry essentielles}

\begin{lstlisting}[style=bash]
# Compiler les contrats
forge build

# Lancer tous les tests
forge test -vv

# Tester seulement un contrat specifique
forge test --match-contract TestDilithium -vv

# Tester une seule fonction
forge test --match-test test_verify_valid_signature -vvv

# Rapport gas detaille
forge test --gas-report

# Profil lite (plus rapide)
FOUNDRY_PROFILE=lite forge test -vv
\end{lstlisting}

%══════════════════════════════════════════════════════════════════════
\section{Exercices}
%══════════════════════════════════════════════════════════════════════

\begin{exercice}{1 --- Lecture de contrat Solidity}
Ouvre \texttt{src/ZKNOX\_dilithium\_utils.sol}.
\begin{enumerate}[label=\alph*)]
  \item Identifie toutes les constantes Dilithium définies.
        Vérifie qu'elles correspondent aux valeurs du Chapitre 5.
  \item Quelle est la structure \texttt{PubKey} ? Quels champs contient-elle ?
  \item Quelle est la structure \texttt{Signature} ? Quels champs contient-elle ?
  \item Trouve la fonction \texttt{VECMULMOD} --- que fait-elle ?
\end{enumerate}
\end{exercice}

\begin{exercice}{2 --- Coûts gas}
\begin{enumerate}[label=\alph*)]
  \item Calcule le coût en ETH d'une vérification ETHDilithium (6.6M gas)
        si le gas price est 20 Gwei et 1 ETH = 3000\$.
  \item Même calcul pour Falcon (350K gas).
  \item Sur un L2 comme Arbitrum (gas 10x moins cher), quel serait le coût ?
\end{enumerate}
\end{exercice}

\begin{solution}
\begin{lstlisting}[style=python]
# Parametres
gas_ethdilithium = 6_600_000
gas_falcon = 350_000
gas_price_gwei = 20
eth_usd = 3000

# ETHDilithium
cost_eth = gas_ethdilithium * gas_price_gwei * 1e-9  # ETH
cost_usd = cost_eth * eth_usd
print(f"ETHDilithium : {cost_eth:.4f} ETH = ${cost_usd:.2f}")
# ~ 0.132 ETH = $396

# Falcon
cost_eth_f = gas_falcon * gas_price_gwei * 1e-9
cost_usd_f = cost_eth_f * eth_usd
print(f"Falcon : {cost_eth_f:.4f} ETH = ${cost_usd_f:.2f}")
# ~ 0.007 ETH = $21

# L2 (10x moins cher)
print(f"Falcon L2 : ${cost_usd_f / 10:.2f}")
# ~ $2.10
\end{lstlisting}
\end{solution}

\begin{exercice}{3 --- ERC-4337}
\begin{enumerate}[label=\alph*)]
  \item Quel est le rôle du \texttt{Bundler} dans ERC-4337 ?
  \item Pourquoi le Smart Wallet paie-t-il le gas (via le Paymaster)
        plutôt que l'utilisateur directement ?
  \item Pour SuVa, quel est l'avantage de ERC-4337 vs modifier le
        protocole Ethereum directement ?
\end{enumerate}
\end{exercice}

\begin{exercice}{4 --- Installer et compiler le projet}
\begin{enumerate}[label=\alph*)]
  \item Installe Foundry et compile les contrats : \texttt{forge build}.
  \item Y a-t-il des erreurs de compilation ? Si oui, lesquelles ?
  \item Lance \texttt{FOUNDRY\_PROFILE=lite forge test -vv} et observe la sortie.
\end{enumerate}
\end{exercice}

\end{document}

\documentclass[12pt, a4paper]{article}
\usepackage[utf8]{inputenc}
\usepackage[T1]{fontenc}
\usepackage[french]{babel}
\usepackage[margin=2.5cm]{geometry}
\usepackage{amsmath, amssymb, amsthm, mathtools}
\usepackage{listings}
\usepackage{xcolor}
\usepackage{tcolorbox}
\tcbuselibrary{theorems, skins, breakable}
\usepackage{booktabs}
\usepackage{fancyhdr}
\usepackage{hyperref}
\usepackage{tikz}
\usetikzlibrary{arrows.meta, positioning, shapes.geometric}

\definecolor{suva_blue}{RGB}{30, 100, 180}
\definecolor{code_bg}{RGB}{245, 247, 250}
\definecolor{code_kw}{RGB}{0, 80, 200}
\definecolor{code_str}{RGB}{180, 40, 40}
\definecolor{code_cmt}{RGB}{80, 130, 80}
\definecolor{warn_bg}{RGB}{255, 248, 220}
\definecolor{success_bg}{RGB}{230, 255, 235}
\definecolor{falcon_color}{RGB}{180, 40, 120}
\definecolor{falcon_bg}{RGB}{255, 240, 250}

\lstdefinestyle{python}{
  language=Python,
  basicstyle=\ttfamily\small,
  keywordstyle=\color{code_kw}\bfseries,
  stringstyle=\color{code_str},
  commentstyle=\color{code_cmt}\itshape,
  backgroundcolor=\color{code_bg},
  frame=single, rulecolor=\color{gray!40},
  numbers=left, numberstyle=\tiny\color{gray},
  breaklines=true, tabsize=4, showstringspaces=false,
}
\lstdefinestyle{output}{
  basicstyle=\ttfamily\small,
  backgroundcolor=\color{gray!10},
  frame=single, rulecolor=\color{gray!50},
  breaklines=true,
}

\newtcolorbox{defbox}[2][]{colback=success_bg,colframe=green!50!black,
  fonttitle=\bfseries,title={Définition : #2},#1}
\newtcolorbox{falconbox}[2][]{colback=falcon_bg,colframe=falcon_color,
  fonttitle=\bfseries,title={Falcon : #2},#1}
\newtcolorbox{exercice}[2][]{colback=white,colframe=suva_blue!60,
  fonttitle=\bfseries,title={Exercice #2},#1,breakable}
\newtcolorbox{solution}[1][]{colback=gray!5,colframe=gray!40,
  fonttitle=\bfseries\itshape,title={Solution},#1,breakable}
\newtcolorbox{importantbox}[1][]{colback=suva_blue!10,colframe=suva_blue,
  fonttitle=\bfseries,title={Connexion avec SuVa},#1}
\newtcolorbox{warningbox}[1][]{colback=warn_bg,colframe=orange!80!black,
  fonttitle=\bfseries,title={Attention},#1}
\newtcolorbox{compbox}[1][]{colback=gray!5,colframe=gray!60,
  fonttitle=\bfseries,title={Comparaison Falcon vs Dilithium},#1}

\pagestyle{fancy}
\fancyhf{}
\fancyhead[L]{\color{suva_blue}\textbf{SuVa --- Chapitre 7}}
\fancyhead[R]{\color{gray}Falcon / FN-DSA}
\fancyfoot[L]{\footnotesize\color{gray}Mohamed Lamine MBENGUE}
\fancyfoot[C]{\thepage}
\fancyfoot[R]{\footnotesize\color{gray}portfolio-alamine.web.app}
\renewcommand{\headrulewidth}{0.4pt}
\renewcommand{\footrulewidth}{0.3pt}

% ══════════════════════════════════════════════════════════════════════
\begin{document}

\begin{center}
  {\color{suva_blue}\rule{\linewidth}{2pt}}\\[0.5cm]
  {\Huge\bfseries\color{suva_blue} Chapitre 7}\\[0.3cm]
  {\large\bfseries Falcon / FN-DSA (NIST FIPS 206)}\\[0.2cm]
  {\small\color{gray} La prochaine étape du projet SuVa --- la cible production}\\[0.3cm]
  {\color{suva_blue}\rule{\linewidth}{2pt}}
\end{center}

\vspace{0.3cm}
\begin{tcolorbox}[colback=suva_blue!5, colframe=suva_blue!40,
  left=6pt, right=6pt, top=4pt, bottom=4pt]
\small
\begin{tabular}{ll}
  \textbf{Auteur}      & Mohamed Lamine MBENGUE \\
  \textbf{Spécialités} & Sécurité · DevOps · Dev Full Stack Web \& Mobile \\
  \textbf{Contact}     & +221 78 178 44 36 · mbenguemohamedlamine2022@gmail.com \\
  \textbf{Portfolio}   & \href{https://portfolio-alamine.web.app}{\texttt{portfolio-alamine.web.app}} \\
\end{tabular}
\end{tcolorbox}
\vspace{0.3cm}

\tableofcontents
\newpage

%══════════════════════════════════════════════════════════════════════
\section{Pourquoi Falcon est la cible de SuVa}
%══════════════════════════════════════════════════════════════════════

\begin{center}
\begin{tabular}{lrrrll}
\toprule
\textbf{Schéma} & \textbf{Gas} & \textbf{Sig (bytes)} & \textbf{pk (bytes)}
  & \textbf{NIST} & \textbf{Sécurité} \\
\midrule
ECDSA (actuel) & 3K & 65 & 33 & — & Cassable quantique \\
\textbf{Falcon-512} & \textbf{350K} & \textbf{666} & \textbf{897} & FIPS 206 & Lattice (NTRU) \\
Dilithium2 & 13.5M & 2420 & 1312 & FIPS 204 & Lattice (MLWE) \\
SPHINCS+ & 50--500M & 7856+ & 32 & FIPS 205 & Hash-based \\
\bottomrule
\end{tabular}
\end{center}

\textbf{Falcon est $\sim 20\times$ moins cher que Dilithium et $\sim 116\times$
plus compact en signature.} C'est pourquoi SuVa vise Falcon comme schéma
principal pour les transactions quotidiennes.

%══════════════════════════════════════════════════════════════════════
\section{NTRU --- La base mathématique de Falcon}
%══════════════════════════════════════════════════════════════════════

\subsection{Différence fondamentale avec Dilithium}

\begin{compbox}
\begin{center}
\begin{tabular}{lll}
\toprule
& \textbf{Dilithium} & \textbf{Falcon} \\
\midrule
Anneau & $R_q = \mathbb{Z}_q[x]/(x^n+1)$ & $R_q = \mathbb{Z}_q[x]/(x^n+1)$ \\
Problème dur & MLWE, MSIS & NTRU, SIS \\
Clé publique & $(rho, t_1)$ & $h = g \cdot f^{-1}$ \\
Structure clé secrète & $(s_1, s_2)$ petits vecteurs & $(f, g)$ polynômes courts \\
Keygen & Rapide & Lent (FFT Gram-Schmidt) \\
Signe & Boucle de rejet & Gaussian sampling \\
Signature & $(c, z, h)$ & $(s_1, s_2)$ court \\
\bottomrule
\end{tabular}
\end{center}
\end{compbox}

\subsection{Le problème NTRU}

\begin{defbox}{Problème NTRU}
Dans $R_q = \mathbb{Z}_q[x]/(x^n + 1)$, connaissant~:
$$h = g \cdot f^{-1} \pmod{q}$$
où $f, g$ ont des coefficients petits (Gaussiens), retrouver $f$ ou $g$.

La clé publique est $h$, la clé secrète est $(f, g)$.
\end{defbox}

\begin{importantbox}
\textbf{Paramètres Falcon-512 (niveau NIST I)~:}
\begin{itemize}
  \item $n = 512$, $q = 12\,289$
  \item $f, g$ : coefficients Gaussiens d'écart-type $\sigma \approx 1.17\sqrt{q/2n}$
  \item Sécurité : $\approx 103$ bits (classique)
\end{itemize}

\textbf{Paramètres Falcon-1024 (niveau NIST V)~:}
\begin{itemize}
  \item $n = 1024$, $q = 12\,289$
  \item Sécurité : $\approx 207$ bits (classique)
\end{itemize}
\end{importantbox}

%══════════════════════════════════════════════════════════════════════
\section{KeyGen Falcon --- Génération de la trappe NTRU}
%══════════════════════════════════════════════════════════════════════

\subsection{Vue d'ensemble}

\begin{lstlisting}[style=python]
# KeyGen Falcon (pseudo-code)
def falcon_keygen(n=512, q=12289):
    # 1. Generer f, g avec des coefficients courts (Gaussiens)
    while True:
        f = sample_gaussian_poly(n, sigma=1.17 * (q/(2*n))**0.5)
        g = sample_gaussian_poly(n, sigma=1.17 * (q/(2*n))**0.5)

        # f doit etre inversible dans R_q
        if not is_invertible(f, q, n):
            continue

        # 2. Calculer h = g * f^(-1) mod q
        f_inv = poly_inverse(f, q, n)
        h = poly_mul(g, f_inv, q, n)

        # 3. Generer la trappe NTRU (pour le signing)
        # Calculer (F, G) tels que f*G - g*F = q (equation NTRU)
        F, G = solve_ntru(f, g, q, n)
        # C'est la partie difficile : utilise FFT sur les entiers

        if F is not None and G is not None:
            break

    pk = h          # cle publique : h (897 bytes pour Falcon-512)
    sk = (f, g, F, G)  # cle secrete : la trappe NTRU
    return pk, sk
\end{lstlisting}

\begin{warningbox}
La génération de clés Falcon est \textbf{significativement plus complexe}
que Dilithium car elle nécessite la décomposition LDL (via FFT sur les
entiers de Gauss). C'est pourquoi on ne génère les clés qu'une fois
et on les réutilise.
\end{warningbox}

%══════════════════════════════════════════════════════════════════════
\section{Sign Falcon --- Gaussian Sampling}
%══════════════════════════════════════════════════════════════════════

Falcon utilise le \textbf{Gaussian sampling} (pas de boucle de rejet
comme Dilithium) pour produire des signatures très courtes.

\subsection{Algorithme de signature}

\begin{lstlisting}[style=python]
# Sign Falcon (pseudo-code tres simplifie)
def falcon_sign(sk, message, n=512, q=12289):
    f, g, F, G = sk

    # 1. Hacher le message -> point cible
    salt = random_bytes(40)
    c = hash_to_point(message, salt, n, q)
    # c est un polynome dans R_q

    # 2. Gaussian preimage sampling
    # Trouver (s1, s2) courts tels que s1*h + s2 = c mod q
    # avec (s1, s2) distribues gaussianement
    # On utilise la trappe (f, g, F, G) pour echantillonner

    # La trappe permet de centrer la Gaussienne autour de la
    # "bonne" solution (la plus courte)
    B = [[g, -f], [G, -F]]  # base NTRU
    # Fast Fourier Sampling dans le reseau NTRU
    s1, s2 = ffsampling_fft(c, B, n, q)

    # Verifier la norme
    sig_norm = l2_norm(s1) + l2_norm(s2)
    # La norme doit etre <= 1.1 * beta_NTRU
    # Si trop grande, recommencer (rare)

    # Signature = (salt, s2) -- s1 peut etre retrouve de h et s2
    return bytes(salt) + pack_int12(s2)
\end{lstlisting}

\begin{importantbox}
\textbf{Gaussian Sampling vs Rejection Sampling~:}
\begin{itemize}
  \item Falcon~: la signature suit une distribution Gaussienne autour
        de la solution optimale. Presque toujours valide du premier coup.
  \item Dilithium~: on échantillonne uniformément et on rejette si trop grand.
        Nécessite en moyenne 4.25 itérations.
\end{itemize}
\textbf{Résultat~:} Falcon produit des signatures \textbf{3--4x plus courtes}
(666 vs 2420 bytes) et s'exécute en \textbf{1 itération}.
\end{importantbox}

%══════════════════════════════════════════════════════════════════════
\section{Verify Falcon --- Vérification simple}
%══════════════════════════════════════════════════════════════════════

La vérification Falcon est \textbf{beaucoup plus simple} que la signature~:

\begin{lstlisting}[style=python]
# Verify Falcon
def falcon_verify(pk, message, signature, n=512, q=12289):
    h = pk  # cle publique

    # Decoder la signature
    salt = signature[:40]
    s2 = unpack_int12(signature[40:])

    # Recomputer s1 : s1 = h^(-1) * (c - s2) mod q
    c = hash_to_point(message, salt, n, q)
    s1 = poly_mul(poly_inverse(h, q, n), c - s2, q, n)

    # Verifier la norme : ||(s1, s2)||_2 <= beta
    norm_sq = sum(x*x for x in s1) + sum(x*x for x in s2)
    # beta^2 est fixe dans les parametres
    return norm_sq <= BETA_SQ
\end{lstlisting}

\textbf{C'est la raison du faible coût gas~:} la vérification ne fait
qu'une multiplication polynomiale + vérification de norme.

%══════════════════════════════════════════════════════════════════════
\section{ETHFALCON --- implémentation EVM disponible}
%══════════════════════════════════════════════════════════════════════

\begin{importantbox}
ZKNox a déjà implémenté Falcon pour l'EVM (ETHFALCON, disponible sur
GitHub : \texttt{ZKNoxHQ/ETHFALCON}).

Pour SuVa, la prochaine étape est~:
\begin{enumerate}
  \item Étudier ETHFALCON (Solidity + Python)
  \item L'intégrer dans le smart wallet SuVa (ERC-4337)
  \item Adapter le SDK Python pour signer avec Falcon
  \item Générer les vecteurs de test Python → Solidity
\end{enumerate}
\end{importantbox}

\subsection{Structure anticipée de l'intégration Falcon dans SuVa}

\begin{lstlisting}[style=python]
# Futur SDK SuVa (a construire)

# Cote Python (off-chain)
from suva_sdk import FalconWallet

wallet = FalconWallet()
pk, sk = wallet.generate_keys()     # Falcon-512 keygen
print(f"Cle publique : {len(pk)} bytes")  # 897 bytes

# Signer une transaction
tx = {"to": "0x...", "value": 1000, "nonce": 42}
sig = wallet.sign_transaction(sk, tx)  # Falcon sign
print(f"Signature : {len(sig)} bytes")   # ~666 bytes

# Cote Solidity (on-chain) - contracte SuVa Wallet
# validateUserOp() appelle verify() du contrat ETHFALCON
# ~350K gas
\end{lstlisting}

%══════════════════════════════════════════════════════════════════════
\section{Roadmap Falcon dans SuVa}
%══════════════════════════════════════════════════════════════════════

\begin{center}
\begin{tikzpicture}[
  step/.style={draw=suva_blue, rounded corners, fill=blue!5,
    minimum width=5cm, minimum height=1cm, text width=4.8cm,
    align=left, font=\small, inner sep=6pt},
  done/.style={step, fill=success_bg, draw=green!50!black},
  todo/.style={step, fill=warn_bg, draw=orange!80},
  critical/.style={step, fill=red!10, draw=red!60},
  arrow/.style={->, thick}
]
  \node[done] (a) {
    \textbf{[FAIT]} Dilithium Python\\
    keygen + sign + verify
  };
  \node[done, below=0.6cm of a] (b) {
    \textbf{[FAIT]} Dilithium Solidity (ZKNox)\\
    vérification on-chain
  };
  \node[todo, below=0.6cm of b] (c) {
    \textbf{[TODO]} Intégrer ETHFALCON\\
    dans le smart wallet SuVa
  };
  \node[todo, below=0.6cm of c] (d) {
    \textbf{[TODO]} SDK client Falcon\\
    (Python, puis JS/TypeScript)
  };
  \node[critical, below=0.6cm of d] (e) {
    \textbf{[TODO]} Vault contracts\\
    wrapper USDT/USDC → suUSDT
  };
  \node[critical, below=0.6cm of e] (f) {
    \textbf{[TODO]} Audit de sécurité\\
    avant tout déploiement
  };

  \draw[arrow] (a) -- (b);
  \draw[arrow] (b) -- (c);
  \draw[arrow] (c) -- (d);
  \draw[arrow] (d) -- (e);
  \draw[arrow] (e) -- (f);
\end{tikzpicture}
\end{center}

%══════════════════════════════════════════════════════════════════════
\section{Comparaison finale : Dilithium vs Falcon vs SPHINCS+}
%══════════════════════════════════════════════════════════════════════

\begin{center}
\begin{tabular}{llll}
\toprule
& \textbf{Dilithium/ML-DSA} & \textbf{Falcon/FN-DSA} & \textbf{SPHINCS+} \\
\midrule
\textbf{Sécurité} & MLWE + MSIS & NTRU + SIS & Uniquement hash \\
\textbf{Hypothèse} & Module-LWE & NTRU & Fonctions de hachage \\
\textbf{Quantum safe} & Oui & Oui & Oui (le plus sûr) \\
\textbf{Pk size (L1)} & 1312 bytes & 897 bytes & 32 bytes \\
\textbf{Sig size (L1)} & 2420 bytes & 666 bytes & 7856 bytes \\
\textbf{Keygen speed} & Rapide & Lent & Rapide \\
\textbf{Sign speed} & Moyen & Rapide (1 iter) & Lent \\
\textbf{Verify speed} & Moyen & Rapide & Lent \\
\textbf{Gas EVM} & 13.5M & \textbf{350K} & 50--500M \\
\textbf{Rôle SuVa} & Recherche & \textbf{Production} & Fallback \\
\bottomrule
\end{tabular}
\end{center}

%══════════════════════════════════════════════════════════════════════
\section{Exercices}
%══════════════════════════════════════════════════════════════════════

\begin{exercice}{1 --- Vérifier les tailles de Falcon}
Falcon-512 produit des signatures de $\lceil 666 \rceil$ bytes.
\begin{enumerate}[label=\alph*)]
  \item La signature contient un sel de 40 bytes + $s_2$ encodé sur
        $\lceil n \log_2(2 \cdot 6146 + 1) \rceil$ bits.
        Pour $n = 512$, calcule la taille théorique de $s_2$.
  \item Pourquoi la signature Falcon est-elle $3.6\times$ plus petite
        que Dilithium2 malgré que $n$ soit $2\times$ plus grand ?
\end{enumerate}
\end{exercice}

\begin{solution}
\begin{lstlisting}[style=python]
import math

# a) Taille de s2
n = 512
# Les coefficients de s2 sont dans [-6146, 6146] approximativement
# encoding compact sur 12 bits par coeff (~12*512/8 = 768 bytes)
# + sel 40 bytes + compression
bits_s2 = n * 12  # approximation
bytes_s2 = bits_s2 // 8
print(f"s2 approximatif : {bytes_s2} bytes")  # ~768
print(f"+ sel 40 bytes = {bytes_s2 + 40} bytes")

# b) Falcon est plus compact car :
# - Distribution Gaussienne -> coefficients plus petits en moyenne
# - Pas de z (qui necessite 18 bits/coeff dans Dilithium)
# - Pas de hint h
# - Pas de c_tilde 32 bytes (le sel remplace)
# Dilithium : c(32) + z(4*256*18/8=2304) + h(~84) = 2420 bytes
# Falcon-512 : sel(40) + s2(~626) = ~666 bytes
\end{lstlisting}
\end{solution}

\begin{exercice}{2 --- Explorer ETHFALCON}
ETHFALCON est disponible sur \texttt{github.com/ZKNoxHQ/ETHFALCON}.
\begin{enumerate}[label=\alph*)]
  \item Clone le repo et explore sa structure.
  \item Quels sont les contrats Solidity clés ?
  \item Y a-t-il une implémentation Python de référence ?
  \item Comment les tests sont-ils organisés ?
  \item Quel est le gas cost mesuré par les tests ?
\end{enumerate}
\end{exercice}

\begin{exercice}{3 --- Plan d'intégration}
Conçois le plan d'intégration de Falcon dans SuVa~:
\begin{enumerate}[label=\alph*)]
  \item Quels fichiers Python du projet actuel devront être adaptés
        pour Falcon ?
  \item Quels contrats Solidity de ETHDILITHIUM sont réutilisables pour
        Falcon (NTT, utils...) ?
  \item Comment les test vectors Python → Solidity seront-ils générés
        pour Falcon ?
  \item Quelles modifications faudra-t-il apporter au smart wallet
        ERC-4337 ?
\end{enumerate}
\end{exercice}

\begin{exercice}{4 --- SPHINCS+ comme fallback}
SPHINCS+ est le fallback "nuclear bunker" de SuVa.
Sa sécurité repose uniquement sur les fonctions de hachage.
\begin{enumerate}[label=\alph*)]
  \item Pourquoi SPHINCS+ est-il le plus sûr des trois schémas ?
  \item Pourquoi son coût gas est-il si élevé (50--500M) ?
  \item Dans quel scénario SuVa utiliserait-il SPHINCS+ plutôt que Falcon ?
  \item SPHINCS+ a des variantes "small" (petite signature, lent) et
        "fast" (grande signature, rapide). Laquelle préférer pour SuVa ?
\end{enumerate}
\end{exercice}

\end{document}


\end{document}
